% Options for packages loaded elsewhere
\PassOptionsToPackage{unicode}{hyperref}
\PassOptionsToPackage{hyphens}{url}
%
\documentclass[
  11pt,a4paper,UTF8,fontset=fandol]{ctexart}
\usepackage{amsmath,amssymb}
\usepackage{setspace}
\usepackage{iftex}
\ifPDFTeX
  \usepackage[T1]{fontenc}
  \usepackage[utf8]{inputenc}
  \usepackage{textcomp} % provide euro and other symbols
\else % if luatex or xetex
  \usepackage{unicode-math} % this also loads fontspec
  \defaultfontfeatures{Scale=MatchLowercase}
  \defaultfontfeatures[\rmfamily]{Ligatures=TeX,Scale=1}
\fi
\usepackage{lmodern}
\ifPDFTeX\else
  % xetex/luatex font selection
\fi
% Use upquote if available, for straight quotes in verbatim environments
\IfFileExists{upquote.sty}{\usepackage{upquote}}{}
\IfFileExists{microtype.sty}{% use microtype if available
  \usepackage[]{microtype}
  \UseMicrotypeSet[protrusion]{basicmath} % disable protrusion for tt fonts
}{}
\makeatletter
\@ifundefined{KOMAClassName}{% if non-KOMA class
  \IfFileExists{parskip.sty}{%
    \usepackage{parskip}
  }{% else
    \setlength{\parindent}{0pt}
    \setlength{\parskip}{6pt plus 2pt minus 1pt}}
}{% if KOMA class
  \KOMAoptions{parskip=half}}
\makeatother
\usepackage{xcolor}
\usepackage[top=24mm,bottom=24mm,left=24mm,right=24mm]{geometry}
\usepackage{longtable,booktabs,array}
\usepackage{calc} % for calculating minipage widths
% Correct order of tables after \paragraph or \subparagraph
\usepackage{etoolbox}
\makeatletter
\patchcmd\longtable{\par}{\if@noskipsec\mbox{}\fi\par}{}{}
\makeatother
% Allow footnotes in longtable head/foot
\IfFileExists{footnotehyper.sty}{\usepackage{footnotehyper}}{\usepackage{footnote}}
\makesavenoteenv{longtable}
\setlength{\emergencystretch}{3em} % prevent overfull lines
\providecommand{\tightlist}{%
  \setlength{\itemsep}{0pt}\setlength{\parskip}{0pt}}
\setcounter{secnumdepth}{-\maxdimen} % remove section numbering
\ifLuaTeX
\usepackage[bidi=basic]{babel}
\else
\usepackage[bidi=default]{babel}
\fi
\babelprovide[main,import]{chinese}
% get rid of language-specific shorthands (see #6817):
\let\LanguageShortHands\languageshorthands
\def\languageshorthands#1{}
\usepackage{fancyhdr}
\usepackage{needspace}
\pagestyle{fancy}
\fancyhf{}
\fancyhead[L]{\small 随机过程五书逐题详解}
\fancyhead[R]{\small 第004批：Durrett 1.56—1.77}
\fancyfoot[C]{\thepage}
\setlength{\headheight}{15pt}
\setlength{\emergencystretch}{3em}
\allowdisplaybreaks[2]
\setlength{\tabcolsep}{4pt}
\renewcommand{\arraystretch}{1.12}
\ctexset{section={format=\Large\bfseries,beforeskip=1.6ex,afterskip=1.0ex},subsection={format=\large\bfseries,beforeskip=1.2ex,afterskip=.7ex}}
\AddToHook{cmd/section/before}{\Needspace{8\baselineskip}}
\AddToHook{cmd/subsection/before}{\Needspace{5\baselineskip}}
\renewcommand{\contentsname}{题目与内容索引}
\ifLuaTeX
  \usepackage{selnolig}  % disable illegal ligatures
\fi
\usepackage{bookmark}
\IfFileExists{xurl.sty}{\usepackage{xurl}}{} % add URL line breaks if available
\urlstyle{same}
\hypersetup{
  pdftitle={随机过程五书逐题详解},
  pdfauthor={中文学习讲义},
  pdflang={zh-CN},
  hidelinks,
  pdfcreator={LaTeX via pandoc}}

\title{随机过程五书逐题详解}
\usepackage{etoolbox}
\makeatletter
\providecommand{\subtitle}[1]{% add subtitle to \maketitle
  \apptocmd{\@title}{\par {\large #1 \par}}{}{}
}
\makeatother
\subtitle{第004批：Durrett第1章收尾 · 1.56---1.77}
\author{中文学习讲义}
\date{22道原题全部小问 · 66道新变式 · 补齐前7题的21道变式}

\begin{document}
\maketitle

{
\setcounter{tocdepth}{1}
\tableofcontents
}
\setstretch{1.12}
\section{本批范围与版本说明}\label{ux672cux6279ux8303ux56f4ux4e0eux7248ux672cux8bf4ux660e}

本批继续 Richard Durrett《Essentials of Stochastic
Processes》第三版（Springer，2016，ISBN
9783319456133）。新增原题为1.56---1.77，共22个题号，包含各题全部小问。每题附三道自编变式及解答，共66道；末尾另补1.1---1.7的21道变式。新增训练合计87道，不重复计入原题数量。

第一章最后一个习题题号是1.77。本批与前三批共同覆盖1.1---1.77，并使第一章77个题号各有三道变式，合计231道。这里的``覆盖''包括对不完整条件或题面错误作明确说明，并不把错误命题硬证明成正确命题。前批1.9(c)的非法转移矩阵仍保留其原样无解结论。

\subsection{原题定位与所用版本}\label{ux539fux9898ux5b9aux4f4dux4e0eux6240ux7528ux7248ux672c}

\begin{longtable}[]{@{}lll@{}}
\toprule\noalign{}
2016版原题号 & 2016版印刷页 & 2021作者公开稿印刷页 \\
\midrule\noalign{}
\endhead
\bottomrule\noalign{}
\endlastfoot
1.56---1.57 & 90 & 73---74 \\
1.58 & 90---91 & 74 \\
1.59---1.62 & 91 & 74---75 \\
1.63---1.67 & 92 & 75 \\
1.68 & 92---93 & 76 \\
1.69---1.71 & 93 & 76 \\
1.72 & 94 & 76 \\
1.73---1.77 & 94 & 77 \\
\end{longtable}

本批以2016年版可读文本核对题号、条件与小问，并查看作者公开稿相应页面图像，核对矩阵、乘积指标和公式。公开稿的页码不能直接当成2016年版页码。文中只简述求解所需条件，主体为独立推导，不复制现成习题答案。

来源：{[}S1{]}
Durrett，2016第三版，第1章第1.13节，印刷页90---94；其可读文本载于
\url{https://dokumen.pub/essentials-of-stochastic-processes-3ed-331945613x-978-3-319-45613-3-978-3-319-45614-0-237-241-244-2.html}。{[}S2{]}
作者公开稿 Version 3.9, May
2021，印刷页73---77及相应正文：\url{https://sites.math.duke.edu/~rtd/EOSP/EOSP2021.pdf}。以下每题的``书内依据''指相同主题所在的原书章节；不以公开稿中错置的例题交叉引用代替实际主题核对。

\subsection{本批特别核对的条件}\label{ux672cux6279ux7279ux522bux6838ux5bf9ux7684ux6761ux4ef6}

1.64中的突变参数属于概率。通常取\(0<u,v<1\)；若只写\(u,v>0\)却允许\(u=v=1\)，原题的收敛断言有反例。1.65的均值收敛要排除\(N=1\)的一球交替情形。1.70中的常返判据使用不可约的生灭链条件。1.71的``其他情况暂留''以\(c>0\)为前提；\(c=0\)另行说明。1.75的常返结论先针对包含0的可达闭类，不把不可达状态也说成常返。

\section{预备工具：本批反复使用的五件事}\label{ux9884ux5907ux5de5ux5177ux672cux6279ux53cdux590dux4f7fux7528ux7684ux4e94ux4ef6ux4e8b}

\subsection{一、到达与正返回要分开}\label{ux4e00ux5230ux8fbeux4e0eux6b63ux8fd4ux56deux8981ux5206ux5f00}

定义

\[
V_A=\inf\{n\ge0:X_n\in A\},\qquad T_x=\inf\{n\ge1:X_n=x\}.
\]

已经从\(x\)出发时，\(V_{\{x\}}=0\)；但\(T_x\)要求至少走一步。概率为1的结论允许零概率例外；无穷大的平均时间也不是``肯定永不返回''。

若\(\mathbb P_x(T_x<\infty)=1\)，称\(x\)常返。若再有\(\mathbb E_xT_x<\infty\)，称正常返；若必返而平均返回时间无穷大，称零常返。

\subsection{二、第一步方程为什么能用}\label{ux4e8cux7b2cux4e00ux6b65ux65b9ux7a0bux4e3aux4ec0ux4e48ux80fdux7528}

原书第1.9节定理1.28处理先到哪个集合，第1.10节定理1.29处理等待多久。令\(A,B\)为停止集合，\(C\)为其余状态。设

\[
h_i=\mathbb P_i(V_A<V_B),\qquad g_i=\mathbb E_iV_{A\cup B}.
\]

边界是\(h|_A=1,h|_B=0,g|_{A\cup B}=0\)。按下一步分类得

\[
h_i=\sum_jp_{ij}h_j,\qquad g_i=1+\sum_jp_{ij}g_j\quad(i\in C).
\]

时间方程多一个1，因为这一步已经用掉了。概率方程没有这个1。

若\(C\)有限，并且从每个状态都存在一条正概率路径到停止集合，可取共同的有限步数\(L\)和正数\(\delta\)，使每\(L\)步内结束的条件概率至少为\(\delta\)。于是

\[
\mathbb P_i(V_{A\cup B}>kL)\le(1-\delta)^k.
\]

这既说明最终一定结束，也说明平均时间有限。若\(Q\)是\(C\)内的转移子矩阵，则

\[
F=(I-Q)^{-1}=I+Q+Q^2+\cdots.
\]

其元素\(F_{ij}\)是结束前在\(j\)的平均停留次数，包含时刻0。概率、时间和累计费用分别为\(Fr,F\mathbf1,Fc\)。这是线性方程的简写，不要求读者先会抽象矩阵理论。

\subsection{三、无向图上不是每个点等可能}\label{ux4e09ux65e0ux5411ux56feux4e0aux4e0dux662fux6bcfux4e2aux70b9ux7b49ux53efux80fd}

在有限连通无向图中，点\(x\)有\(d(x)\)个邻居，每次等概率选一个邻居，则

\[
\pi(x)=\frac{d(x)}{D},\qquad D=\sum_zd(z).
\]

因为相邻\(x,y\)满足

\[
\pi(x)p(x,y)=\frac{d(x)}D\frac1{d(x)}=\frac1D=\pi(y)p(y,x).
\]

对\(x\)求和即得平稳方程。这是第1.5节细致平衡的应用。由第1.6---1.7节定理1.21，正返回平均时间为

\[
\mathbb E_xT_x=\frac1{\pi(x)}.
\]

这个公式不要求非周期。非周期性是全序列分布收敛所需的另一项条件。

\subsection{四、几何等待与尾概率求和}\label{ux56dbux51e0ux4f55ux7b49ux5f85ux4e0eux5c3eux6982ux7387ux6c42ux548c}

若\(G\)取\(1,2,\ldots\)，每次成功概率为\(s>0\)，则

\[
\mathbb P(G=k)=(1-s)^{k-1}s,\quad
\mathbb EG=\frac1s,\quad
\operatorname{Var}(G)=\frac{1-s}{s^2}.
\]

由等比级数求导可以逐项得到后两个公式。对于任意非负整数随机变量\(T\)，有

\[
T=\sum_{n=0}^{\infty}\mathbf1_{\{T>n\}},\qquad
\mathbb ET=\sum_{n=0}^{\infty}\mathbb P(T>n).
\]

先对有限个非负项求期望，再让项数增加，就得到该式；允许结果是无穷大。

\subsection{五、灭绝概率为什么取最小不动点}\label{ux4e94ux706dux7eddux6982ux7387ux4e3aux4ec0ux4e48ux53d6ux6700ux5c0fux4e0dux52a8ux70b9}

第1.11节例1.55及引理1.31讨论分枝过程。每个个体独立产生\(K\)个后代，令

\[
f(s)=\mathbb E(s^K)=\sum_{k\ge0}p_ks^k.
\]

从一个个体出发，第\(n\)代已经灭绝的概率记为\(q_n\)。有\(q_0=0\)，且

\[
q_{n+1}=f(q_n).
\]

这是因为第一代有\(k\)个个体时，\(k\)条后代支系必须全部灭绝，条件概率为\(q_n^k\)。序列\(q_n\)递增且不超过1，所以有极限\(q\)，并满足\(f(q)=q\)。任意另一个不动点\(r\in[0,1]\)均满足\(q_0\le r\)；由\(f\)递增，归纳得\(q_n\le r\)，故\(q\le r\)。因此必须选\([0,1]\)内最小的根，而不是见到根1就停止。

\section{习题1.56：三种贷款的终止时间与清偿概率}\label{ux4e60ux98981.56ux4e09ux79cdux8d37ux6b3eux7684ux7ec8ux6b62ux65f6ux95f4ux4e0eux6e05ux507fux6982ux7387}

\textbf{书内依据：第1.9节定理1.28、第1.10节定理1.29。原题位置见{[}S1{]}页90、{[}S2{]}页73。}

按\(V,30,15,P,F\)排列，求解所需矩阵为

\[
P=\begin{pmatrix}
.55&.35&0&.05&.05\\
.15&.54&.25&.05&.01\\
.20&0&.75&.04&.01\\
0&0&0&1&0\\
0&0&0&0&1
\end{pmatrix}.
\]

\(P,F\)是两个终止状态；前三个标签是贷款类型，不是``一步等于多少年''。原题未指定每次转移的日历长度，时间答案用``转移期''。

\subsection{（a）三个起点的平均终止时间}\label{aux4e09ux4e2aux8d77ux70b9ux7684ux5e73ux5747ux7ec8ux6b62ux65f6ux95f4}

令\((a,b,c)=(g_V,g_{30},g_{15})\)。第一步方程为

\[
\begin{aligned}
a&=1+.55a+.35b,\\
b&=1+.15a+.54b+.25c,\\
c&=1+.20a+.75c.
\end{aligned}
\]

分别整理为

\[
9a-7b=20,\qquad -15a+46b-25c=100,\qquad -4a+5c=20.
\]

第三式给\(c=4+4a/5\)。代入中间一式得到\(-35a+46b=200\)。再用\(b=(9a-20)/7\)，得到

\[
-245a+414a-920=1400,\qquad169a=2320.
\]

依次代回，得到

\[
\boxed{(g_V,g_{30},g_{15})=\left(\frac{2320}{169},\frac{2500}{169},\frac{2532}{169}\right).}
\]

\subsection{（b）三个起点的最终清偿概率}\label{bux4e09ux4e2aux8d77ux70b9ux7684ux6700ux7ec8ux6e05ux507fux6982ux7387}

现在令\((a,b,c)=(h_V,h_{30},h_{15})\)，边界\(h_P=1,h_F=0\)。方程是

\[
\begin{aligned}
a&=.55a+.35b+.05,\\
b&=.15a+.54b+.25c+.05,\\
c&=.20a+.75c+.04.
\end{aligned}
\]

即

\[
9a-7b=1,\qquad -15a+46b-25c=5,\qquad -20a+25c=4.
\]

由\(c=4/25+4a/5\)得\(-35a+46b=9\)。代入\(b=(9a-1)/7\)，有\(169a=109\)。因此

\[
\boxed{(h_V,h_{30},h_{15})=\left(\frac{109}{169},\frac{116}{169},\frac{2856}{4225}\right).}
\]

三种起点的最终止赎概率分别是\(60/169,53/169,1369/4225\)。各自与清偿概率之和为1。

\subsection{变式1.56A：每期费用不同}\label{ux53d8ux5f0f1.56aux6bcfux671fux8d39ux7528ux4e0dux540c}

\textbf{题目。}在三个非终止状态每期费用依次为\(1,2,3\)，终止后不再收费。求从各状态开始的总费用期望，时刻0所在状态也收费。

\textbf{解答。}设费用向量为\(c=(1,2,3)^{\mathsf T}\)，未知期望\(w\)满足\(w=c+Qw\)。原题的\(Q\)给出

\[
F=\frac1{169}\begin{pmatrix}920&700&700\\700&900&900\\736&560&1236\end{pmatrix}.
\]

因此

\[
w=Fc=\frac1{169}\begin{pmatrix}920+1400+2100\\700+1800+2700\\736+1120+3708\end{pmatrix}
=\boxed{\left(\frac{340}{13},\frac{400}{13},\frac{428}{13}\right)^{\mathsf T}}.
\]

不能用``平均费用''随意乘原题的平均时间，因为各状态出现的次数不同。

\subsection{变式1.56B：终止时固定比例清偿}\label{ux53d8ux5f0f1.56bux7ec8ux6b62ux65f6ux56faux5b9aux6bd4ux4f8bux6e05ux507f}

\textbf{题目。}保持\(Q\)不变，但每个状态一旦离开非终止区域，都以同一概率\(\alpha\)清偿，以\(1-\alpha\)止赎。求清偿概率。

\textbf{解答。}令\(r=\mathbf1-Q\mathbf1\)为一步终止概率列向量。清偿概率满足\(h=Qh+\alpha r\)。代入\(h=\alpha\mathbf1\)，右端为

\[
\alpha Q\mathbf1+\alpha(\mathbf1-Q\mathbf1)=\alpha\mathbf1.
\]

由有限吸收方程唯一性，所有起点清偿概率都是\(\boxed\alpha\)。这说明终止时间可以随起点改变，而结局概率保持相同。

\subsection{变式1.56C：已知最终清偿，平均还要多久}\label{ux53d8ux5f0f1.56cux5df2ux77e5ux6700ux7ec8ux6e05ux507fux5e73ux5747ux8fd8ux8981ux591aux4e45}

\textbf{题目。}从\(V\)开始，条件于最终清偿，求平均终止时间。

\textbf{解答。}不能直接用\(g_V/h_V\)。设\(u_i=\mathbb E_i[T\mathbf1_{\{\text{清偿}\}}]\)。按第一步拆开，已经花掉的一期只有在最终清偿的路径上才计入，故

\[
u=h+Qu,\qquad u=Fh.
\]

代入原题\(h\)和上面的\(F\)得

\[
u_V=\frac{261448}{28561}.
\]

条件期望的定义给

\[
\boxed{\mathbb E_V[T\mid\text{清偿}]=\frac{u_V}{h_V}=\frac{261448}{18421}.}
\]

这里的分母是清偿事件的概率；分子必须先把未清偿路径的时间排除。

\section{习题1.57：理发店链的平稳分布、先到概率与退出时间}\label{ux4e60ux98981.57ux7406ux53d1ux5e97ux94feux7684ux5e73ux7a33ux5206ux5e03ux5148ux5230ux6982ux7387ux4e0eux9000ux51faux65f6ux95f4}

\textbf{书内依据：第1.5节细致平衡，第1.9---1.10节第一步方程。原题{[}S1{]}页90、{[}S2{]}页74。}

按\(0,1,2,3,4\)排列，

\[
P=\begin{pmatrix}
0&1&0&0&0\\3/5&0&2/5&0&0\\0&3/4&0&1/4&0\\0&0&3/4&0&1/4\\0&0&0&3/4&1/4
\end{pmatrix}.
\]

原链中的0和4不是吸收状态。但是，研究``第一次到0或4就停止''时，边界值设在0和4即可。不要为了算退出时间而把修改后的链用于求原链平稳分布。

\subsection{（a）平稳分布}\label{aux5e73ux7a33ux5206ux5e03}

逐对平衡得

\[
\pi_1=\frac53\pi_0,\quad
\pi_2=\frac{(2/5)}{(3/4)}\pi_1=\frac89\pi_0,\quad
\pi_3=\frac8{27}\pi_0,\quad
\pi_4=\frac8{81}\pi_0.
\]

把所有比例通分，非归一化权重为\((81,135,72,24,8)\)，总和320。因此

\[
\boxed{\pi=\frac1{320}(81,135,72,24,8).}
\]

\subsection{（b）先到0而不是4的概率}\label{bux5148ux52300ux800cux4e0dux662f4ux7684ux6982ux7387}

令\(h_0=1,h_4=0\)，内部方程为

\[
h_1=\frac35+\frac25h_2,\quad h_2=\frac34h_1+\frac14h_3,\quad h_3=\frac34h_2.
\]

代入中间一式：

\[
h_2=\frac9{20}+\left(\frac3{10}+\frac3{16}\right)h_2
=\frac9{20}+\frac{39}{80}h_2.
\]

故\((41/80)h_2=9/20\)，再代回得到

\[
\boxed{(h_1,h_2,h_3)=(39/41,36/41,27/41).}
\]

\subsection{（c）到达0或4的平均时间}\label{cux5230ux8fbe0ux62164ux7684ux5e73ux5747ux65f6ux95f4}

边界\(g_0=g_4=0\)，于是

\[
g_1=1+\frac25g_2,\quad g_2=1+\frac34g_1+\frac14g_3,\quad g_3=1+\frac34g_2.
\]

代入第二式得\(g_2=2+(39/80)g_2\)，因此

\[
\boxed{(g_1,g_2,g_3)=(105/41,160/41,161/41).}
\]

\subsection{变式1.57A：不在0停止，直到首次到4}\label{ux53d8ux5f0f1.57aux4e0dux57280ux505cux6b62ux76f4ux5230ux9996ux6b21ux52304}

\textbf{题目。}求原链从0出发首次到4的平均时间。

\textbf{解答。}记\(k_4=0\)。方程\(k_0=1+k_1\)与\(k_1=1+(3/5)k_0+(2/5)k_2\)给\(k_1-k_2=4\)。再由\(k_2=1+(3/4)k_1+(1/4)k_3\)得\(k_2-k_3=16\)。最后

\[
k_3=1+\frac34k_2=1+\frac34(k_3+16)
\]

给\(k_3=52\)。依次有\(k_2=68,k_1=72,k_0=73\)。所以答案为\(\boxed{73}\)，远大于``碰到0或4就停止''的时间。

\subsection{变式1.57B：从0出发的正返回时间}\label{ux53d8ux5f0f1.57bux4ece0ux51faux53d1ux7684ux6b63ux8fd4ux56deux65f6ux95f4}

\textbf{题目。}求\(\mathbb E_0T_0\)，其中\(T_0=\inf\{n\ge1:X_n=0\}\)。

\textbf{解答。}原链有限且不可约。由定理1.21，

\[
\boxed{\mathbb E_0T_0=1/\pi_0=320/81.}
\]

若写成允许\(n=0\)的首次到达时间，则答案会是0；这不是本题的正返回时间。

\subsection{变式1.57C：加入停留，不改变移动方向}\label{ux53d8ux5f0f1.57cux52a0ux5165ux505cux7559ux4e0dux6539ux53d8ux79fbux52a8ux65b9ux5411}

\textbf{题目。}每一步以\(1-\varepsilon\)保持不动，以\(\varepsilon\in(0,1]\)执行一次原转移。三个问题的答案怎样变化？

\textbf{解答。}新矩阵\(P_\varepsilon=(1-\varepsilon)I+\varepsilon P\)，所以\(\pi P_\varepsilon=\pi\)，平稳分布不变。先到概率的方程

\[
h=(1-\varepsilon)h+\varepsilon Ph
\]

约去停留项后与原题相同，所以先到概率不变。退出时间满足\(g'=1+(1-\varepsilon)g'+\varepsilon Qg'\)，即\((I-Q)g'=\mathbf1/\varepsilon\)。因此\(\boxed{g'=g/\varepsilon}\)。

\section{习题1.58：二状态链的精确收敛公式}\label{ux4e60ux98981.58ux4e8cux72b6ux6001ux94feux7684ux7cbeux786eux6536ux655bux516cux5f0f}

\textbf{书内依据：第1.1节马尔可夫性、第1.2节分布递推、第1.4节二状态平稳分布。原题{[}S1{]}页90---91、{[}S2{]}页74。}

令

\[
P=\begin{pmatrix}1-a&a\\b&1-b\end{pmatrix},\quad 0\le a,b\le1,
\qquad u_n=\mathbb P(X_n=1).
\]

按当前状态分类，

\[
u_{n+1}=(1-a)u_n+b(1-u_n)=b+(1-a-b)u_n.
\]

当\(a+b>0\)时，令\(u_*=b/(a+b)\)。直接计算\(b+(1-a-b)u_*=u_*\)，因此

\[
\boxed{u_{n+1}-u_*=(1-a-b)(u_n-u_*).}
\]

现在作归纳：\(n=0\)时公式显然成立；若第\(n\)步成立，乘\(1-a-b\)即得第\(n+1\)步。因此

\[
\boxed{u_n=\frac b{a+b}+(1-a-b)^n\left(u_0-\frac b{a+b}\right).}
\]

若\(0<a+b<2\)，则\(|1-a-b|<1\)，误差以几何速度趋于0。\(a+b=1\)时一步即到平衡。\(a=b=0\)时原状态永不变，不能除以\(a+b\)；\(a=b=1\)时确定性交替，通常不收敛。\(n=0\)的幂取1，即使底数为0也按公式初始项理解。

\subsection{变式1.58A：达到指定误差需要几步}\label{ux53d8ux5f0f1.58aux8fbeux5230ux6307ux5b9aux8befux5deeux9700ux8981ux51e0ux6b65}

\textbf{题目。}设\(0<r=|1-a-b|<1\)、初始误差\(c=|u_0-u_*|>\delta>0\)，求使\(|u_n-u_*|\le\delta\)的最小整数\(n\)。

\textbf{解答。}由原题等式，要求\(cr^n\le\delta\)。取对数，因\(\log r<0\)，除法时不等号改变方向：

\[
n\ge\frac{\log(\delta/c)}{\log r}.
\]

故\(\boxed{n_{\min}=\lceil\log(\delta/c)/\log r\rceil}\)。若\(c\le\delta\)则\(n_{\min}=0\)；若\(r=0,c>\delta\)则\(n_{\min}=1\)。在两状态空间，总变差距离也恰等于\(|u_n-u_*|\)。

\subsection{变式1.58B：前后两次观察的相关性}\label{ux53d8ux5f0f1.58bux524dux540eux4e24ux6b21ux89c2ux5bdfux7684ux76f8ux5173ux6027}

\textbf{题目。}平稳起步，记\(Z_n=\mathbf1_{\{X_n=1\}}\)。求\(\operatorname{Cov}(Z_0,Z_k)\)。

\textbf{解答。}原题条件版本给

\[
\mathbb E(Z_k\mid Z_0)=u_*+(1-a-b)^k(Z_0-u_*).
\]

两边减\(u_*\)，乘\(Z_0-u_*\)后取期望，得到

\[
\boxed{\operatorname{Cov}(Z_0,Z_k)=u_*(1-u_*)(1-a-b)^k.}
\]

\(a+b>1\)时奇数滞后协方差为负，偶数滞后为正；均值收敛并不意味着相邻时刻独立。

\subsection{变式1.58C：由长期比例和振荡速度反求矩阵}\label{ux53d8ux5f0f1.58cux7531ux957fux671fux6bd4ux4f8bux548cux632fux8361ux901fux5ea6ux53cdux6c42ux77e9ux9635}

\textbf{题目。}长期状态1比例为\(2/5\)，每一步误差乘以\(-1/4\)，求\(P\)。

\textbf{解答。}\(1-a-b=-1/4\)给\(a+b=5/4\)；\(b/(a+b)=2/5\)给\(b=1/2\)，从而\(a=3/4\)。因此

\[
\boxed{P=\begin{pmatrix}1/4&3/4\\1/2&1/2\end{pmatrix}.}
\]

两行和为1，所有项均非负，所以它确实是合法矩阵。

\section{习题1.59：一般两盒交换模型}\label{ux4e60ux98981.59ux4e00ux822cux4e24ux76d2ux4ea4ux6362ux6a21ux578b}

\textbf{书内依据：第1.1节建模、第1.5节细致平衡、附录A.1的等可能计数。原题{[}S1{]}页91、{[}S2{]}页74。}

两盒各装\(m\ge1\)个球，总共\(b\)个黑球、\(2m-b\)个白球，其中\(0\le b\le2m\)。状态\(i\)为第一盒黑球数。正确状态空间不是一律\(0,\ldots,m\)，而是

\[
\boxed{\max(0,b-m)\le i\le\min(b,m).}
\]

\subsection{（a）转移概率}\label{aux8f6cux79fbux6982ux7387}

第一盒有\(i\)黑、\(m-i\)白；第二盒有\(b-i\)黑、\(m-b+i\)白。因此

\[
\begin{aligned}
p(i,i+1)&=\frac{(m-i)(b-i)}{m^2},\\
p(i,i-1)&=\frac{i(m-b+i)}{m^2},\\
p(i,i)&=\frac{i(b-i)+(m-i)(m-b+i)}{m^2}.
\end{aligned}
\]

增加来自``第一盒取白、第二盒取黑''，减少来自相反颜色方向，不变来自同色交换。四种颜色组合穷尽全部可能，故每行概率和为1；不可到达的边界外状态的概率为0。

\subsection{（b）验证平稳分布}\label{bux9a8cux8bc1ux5e73ux7a33ux5206ux5e03}

令

\[
\pi_i=\frac{\binom bi\binom{2m-b}{m-i}}{\binom{2m}m}.
\]

把\(2m\)个有编号球分为黑、白两组。选第一盒的\(m\)个球，若其中\(i\)个黑球，可选\(\binom bi\binom{2m-b}{m-i}\)种。把所有合法\(i\)相加恰为\(\binom{2m}m\)，所以\(\sum_i\pi_i=1\)。

相邻项比值为

\[
\frac{\pi_{i+1}}{\pi_i}
=\frac{b-i}{i+1}\frac{m-i}{m-b+i+1}
=\frac{p(i,i+1)}{p(i+1,i)}.
\]

故相邻状态满足细致平衡，其他不同状态之间双向概率都是0，自环自动满足。因此该\(\pi\)平稳。

\subsection{（c）为什么这个计数直觉可信}\label{cux4e3aux4ec0ux4e48ux8fd9ux4e2aux8ba1ux6570ux76f4ux89c9ux53efux4fe1}

微观状态是``第一盒究竟包含哪\(m\)个有编号球''。一次交换可以反向做回来，正向与反向指定交换的概率都是\(1/m^2\)。所以均匀分配的微观状态满足细致平衡。将微观状态按黑球数归组，正好得到上面的超几何概率。这既解释了公式，也说明为什么不能把每个黑球数都当成等可能。

\subsection{变式1.59A：不求完整分布也能求均值}\label{ux53d8ux5f0f1.59aux4e0dux6c42ux5b8cux6574ux5206ux5e03ux4e5fux80fdux6c42ux5747ux503c}

\textbf{题目。}从\(i\)出发，求\(\mu_n=\mathbb E_iX_n\)。

\textbf{解答。}条件平均变化为

\[
\mathbb E(X_{n+1}-X_n\mid X_n=x)
=\frac{(m-x)(b-x)-x(m-b+x)}{m^2}
=\frac bm-\frac{2x}m.
\]

故\(\mu_{n+1}=b/m+(1-2/m)\mu_n\)。减去固定值\(b/2\)并迭代：

\[
\boxed{\mu_n=\frac b2+\left(1-\frac2m\right)^n\left(i-\frac b2\right).}
\]

\(m=1,b=1\)时系数为\(-1\)，均值交替；不要在这一边界声称收敛。

\subsection{变式1.59B：两个盒子的容量不相同}\label{ux53d8ux5f0f1.59bux4e24ux4e2aux76d2ux5b50ux7684ux5bb9ux91cfux4e0dux76f8ux540c}

\textbf{题目。}两盒容量改为\(m,n\ge1\)，总黑球数为\(b\)，仍各抽一个交换。求转移概率和平稳分布。

\textbf{解答。}第一盒\(x\)黑，第二盒\(b-x\)黑，所以

\[
p(x,x+1)=\frac{(m-x)(b-x)}{mn},\quad
p(x,x-1)=\frac{x(n-b+x)}{mn}.
\]

不变概率为1减去这两项。状态范围\(\max(0,b-n)\le x\le\min(m,b)\)。均匀选择第一盒\(m\)个有编号球给

\[
\boxed{\pi_x=\frac{\binom bx\binom{m+n-b}{m-x}}{\binom{m+n}m}.}
\]

相邻比值为\((b-x)(m-x)/[(x+1)(n-b+x+1)]\)，正好等于双向转移之比，因此完整验证了平稳性。

\subsection{变式1.59C：平衡时的方差}\label{ux53d8ux5f0f1.59cux5e73ux8861ux65f6ux7684ux65b9ux5dee}

\textbf{题目。}原题平衡时求\(\operatorname{Var}(X)\)。

\textbf{解答。}给\(b\)个黑球编号，令\(I_j\)表示第\(j\)个黑球是否在第一盒。则\(X=\sum_{j=1}^bI_j\)，且\(\mathbb EI_j=1/2\)。两个不同黑球同时进入第一盒的概率是

\[
\frac m{2m}\frac{m-1}{2m-1}=\frac{m-1}{2(2m-1)}.
\]

故\(\operatorname{Cov}(I_i,I_j)=-1/[4(2m-1)]\)。展开方差：

\[
\operatorname{Var}(X)=\frac b4-\frac{b(b-1)}{4(2m-1)}
=\boxed{\frac{b(2m-b)}{4(2m-1)}}.
\]

这里的负协方差来自盒子容量固定；不能把各黑球的位置误当成独立。

\section{习题1.60：书架``取到就移到最左''的平稳分布}\label{ux4e60ux98981.60ux4e66ux67b6ux53d6ux5230ux5c31ux79fbux5230ux6700ux5de6ux7684ux5e73ux7a33ux5206ux5e03}

\textbf{书内依据：第1.1节状态选择、第1.4节平稳分布；证明另用独立请求与等比级数。原题{[}S1{]}页91、{[}S2{]}页74。}

共有\(n\)本不同书，每次独立请求第\(i\)本的概率为\(p_i>0\)，\(\sum_ip_i=1\)。被请求的书移到最左，其余书保持相对次序。要证明排列\((i_1,\ldots,i_n)\)的平稳概率为

\[
\boxed{\pi(i_1,\ldots,i_n)=\prod_{k=1}^n\frac{p_{i_k}}{1-\sum_{\ell<k}p_{i_\ell}}.}
\]

空和为0，最后一个分母恰是最后一本书的概率，所以最后一个因子为1。

\subsection{第一步：理解书架次序记录了什么}\label{ux7b2cux4e00ux6b65ux7406ux89e3ux4e66ux67b6ux6b21ux5e8fux8bb0ux5f55ux4e86ux4ec0ux4e48}

所有书至少被请求过一次以后，最左边是最近请求的书。删去它的全部较早请求后，下一个最近请求的不同书排第二，如此继续。因此，最终书架等于``从最新请求往过去读，按第一次遇到的不同书排列''。

\subsection{第二步：倒着读独立请求}\label{ux7b2cux4e8cux6b65ux5012ux7740ux8bfbux72ecux7acbux8bf7ux6c42}

设当前已经发现不同书的集合为\(A\)，其请求概率和为\(s_A\)。接下来看到的一串请求可能仍属于\(A\)，直到首次出现新书\(j\notin A\)。首次新书为\(j\)的概率是

\[
p_j+s_Ap_j+s_A^2p_j+\cdots=\frac{p_j}{1-s_A}.
\]

于是第一个不同书为\(i_1\)的概率是\(p_{i_1}\)；第二个为\(i_2\)的条件概率是\(p_{i_2}/(1-p_{i_1})\)；接下来按同样规则相乘，恰得到题目公式。

这些条件概率在每一层对``尚未出现的书''求和都为1，所以乘积公式在全部\(n!\)个排列上自动归一化，不需直接展开\(n!\)项。

\subsection{第三步：为什么这不只是某个概率分布，而确实平稳}\label{ux7b2cux4e09ux6b65ux4e3aux4ec0ux4e48ux8fd9ux4e0dux53eaux662fux67d0ux4e2aux6982ux7387ux5206ux5e03ux800cux786eux5b9eux5e73ux7a33}

从任意固定书架出发，做\(m\)次独立请求，分布记为\(\mu_m\)。某本书\(i\)从未出现的概率为\((1-p_i)^m\)，故至少有一本没出现的概率不超过

\[
\sum_{i=1}^n(1-p_i)^m\longrightarrow0.
\]

把长度\(m\)的请求逆序读，仍是独立同分布。随着\(m\)增加，第一次见全\(n\)种书的时刻几乎必定有限：上面的未见全概率趋于0已经给出了证明。因此倒读所得排列的分布趋于第二步算出的\(\pi\)。在所有书出现后，最终书架不再依赖初始排列，故\(\mu_m\to\pi\)。

又有\(\mu_{m+1}=\mu_mP\)。状态总数\(n!\)有限，可以在每一个有限和中取极限，得到\(\pi=\pi P\)。证明完成。

若有\(p_i=0\)，题目原式可能出现\(0/0\)，不能照抄。正请求概率的书最终全部排在零概率书之前；零概率书之间保持原有次序。对正概率书使用同样公式，再在不同零概率书次序对应的闭类之间任意混合，便得到相应平稳分布。

\subsection{变式1.60A：三本书的六个排列}\label{ux53d8ux5f0f1.60aux4e09ux672cux4e66ux7684ux516dux4e2aux6392ux5217}

\textbf{题目。}\(p_1=1/2,p_2=1/3,p_3=1/6\)，列出平稳分布。

\textbf{解答。}例如\(\pi(1,2,3)=(1/2)(1/3)/(1/2)=1/3\)。逐项同算：

\begin{longtable}[]{@{}lllllll@{}}
\toprule\noalign{}
排列 & 123 & 132 & 213 & 231 & 312 & 321 \\
\midrule\noalign{}
\endhead
\bottomrule\noalign{}
\endlastfoot
平稳概率 & \(1/3\) & \(1/6\) & \(1/4\) & \(1/12\) & \(1/10\) &
\(1/15\) \\
\end{longtable}

通分至60，分子为20、10、15、5、6、4，总和60。六个排列并不等可能。

\subsection{变式1.60B：某一本书的平均位置}\label{ux53d8ux5f0f1.60bux67d0ux4e00ux672cux4e66ux7684ux5e73ux5747ux4f4dux7f6e}

\textbf{题目。}平衡时求书\(i\)从左数的平均位置\(R_i\)。

\textbf{解答。}\(R_i=1+\sum_{j\ne i}\mathbf1_{\{j\text{在}i\text{前}\}}\)。只看\(i,j\)两种请求，最后一次相关请求是\(j\)的概率为\(p_j/(p_i+p_j)\)，这也就是\(j\)在\(i\)前的概率。因此

\[
\boxed{\mathbb ER_i=1+\sum_{j\ne i}\frac{p_j}{p_i+p_j}.}
\]

上题三个平均位置依次为\(33/20,29/15,29/12\)。它们之和为6，等于每个排列的位置和\(1+2+3\)，可作独立检查。

\subsection{变式1.60C：有平稳分布是否就可逆}\label{ux53d8ux5f0f1.60cux6709ux5e73ux7a33ux5206ux5e03ux662fux5426ux5c31ux53efux9006}

\textbf{题目。}三本书请求概率均正，证明此链不满足细致平衡。

\textbf{解答。}\(123\to312\)可以由请求书3实现，概率\(p_3>0\)。但从312出发，一步只可能变成312、132或231，不能回到123。因此

\[
\pi(123)P(123,312)>0,\qquad \pi(312)P(312,123)=0.
\]

两边不相等。故链有平稳分布却不可逆，不能把这两个概念混为一谈。

\section{习题1.61：钟面上的随机游走}\label{ux4e60ux98981.61ux949fux9762ux4e0aux7684ux968fux673aux6e38ux8d70}

\textbf{书内依据：第1.4.1节双随机矩阵，第1.7节正返回时间公式，第1.9节首达概率。}把十二个位置编号为\(0,1,\ldots,11\)，每次以相同概率顺、逆时针走一格，从0开始。

\subsection{（a）第一次正返回的平均时间}\label{aux7b2cux4e00ux6b21ux6b63ux8fd4ux56deux7684ux5e73ux5747ux65f6ux95f4}

这里要计算\(T_0=\min\{n\ge1:X_n=0\}\)，不能用允许\(n=0\)的到达时间。每个位置有两个邻居，转移矩阵的每列和也等于1，故平稳分布为\(\pi_i=1/12\)。有限链不可约，因此

\[
\boxed{\mathbb E_0T_0=\frac1{\pi_0}=12.}
\]

链的周期是2，但正返回均值公式不要求非周期性。不能因为分布逐步振荡就拒绝使用该公式。

\subsection{（b）返回之前走遍另外十一个位置的概率}\label{bux8fd4ux56deux4e4bux524dux8d70ux904dux53e6ux5916ux5341ux4e00ux4e2aux4f4dux7f6eux7684ux6982ux7387}

先假设第一步顺时针到了1。若在首次返回0之前走遍其他位置，就必须先到11，再到0。由于每次只走一格，走到11途中必然经过2至10，不需要另加``中途覆盖''的条件。

在区间\(0,1,\ldots,11\)上记\(h_i\)为先到11再到0的概率，则

\[
h_0=0,\quad h_{11}=1,\quad h_i=\frac12h_{i-1}+\frac12h_{i+1}.
\]

整理得\(h_{i+1}-h_i=h_i-h_{i-1}\)，即相邻差相等。利用端点得到\(h_i=i/11\)，特别是\(h_1=1/11\)。第一步逆时针时，镜面对称给出同样的条件概率。因此

\[
\boxed{\mathbb P_0(\text{正返回前走遍其他位置})
=\tfrac12\cdot\tfrac1{11}+\tfrac12\cdot\tfrac1{11}=\frac1{11}.}
\]

\textbf{易错点：}两种方向各占概率\(1/2\)，不能直接把两个条件概率相加得到\(2/11\)。

\subsection{变式1.61A：任意大小的圆环}\label{ux53d8ux5f0f1.61aux4efbux610fux5927ux5c0fux7684ux5706ux73af}

\textbf{题目。}把钟面改成\(M\ge3\)个位置，仍等概率走向两个邻居，求上述两个量。

\textbf{解答。}列和为1，故\(\pi_i=1/M\)，正返回均值为\(M\)。第一步后，把尚未跨过的起点切开，得到长\(M-1\)的吸收区间。相邻差为常数，起点1先到\(M-1\)的概率为\(1/(M-1)\)。两种方向加权后仍是这一概率。因此答案为

\[\boxed{M,\qquad\frac1{M-1}.}\]

\subsection{变式1.61B：偏向某一个方向}\label{ux53d8ux5f0f1.61bux504fux5411ux67d0ux4e00ux4e2aux65b9ux5411}

\textbf{题目。}顺时针概率\(p\in(0,1)\)，逆时针概率\(q=1-p\)，其他同上。

\textbf{解答。}每列收到的概率仍为\(p+q=1\)，所以返回均值仍是\(M\)。令\(r=q/p\)。在区间\(0,\ldots,M-1\)内，首达概率满足\(p(h_{i+1}-h_i)=q(h_i-h_{i-1})\)，相邻差构成公比\(r\)的等比数列。若\(p\ne q\)，从1出发成功概率为\((1-r)/(1-r^{M-1})\)。第一步逆时针后，左右概率互换。因此总覆盖概率为

\[
p\frac{1-r}{1-r^{M-1}}
+q\frac{1-r^{-1}}{1-r^{-(M-1)}}
=\boxed{\frac{(p-q)(p^{M-1}+q^{M-1})}{p^{M-1}-q^{M-1}}}.
\]

分子、分母在\(p<q\)时都变号，结果仍为正。\(p=q=1/2\)用原题的线性解，概率为\(1/(M-1)\)；不能向等比公式直接代入而保留\(0/0\)。

\subsection{变式1.61C：到对面不是返回}\label{ux53d8ux5f0f1.61cux5230ux5bf9ux9762ux4e0dux662fux8fd4ux56de}

\textbf{题目。}公平圆环上，求从0第一次到距离\(d\)的给定位置的平均时间，\(1\le d\le M-1\)。

\textbf{解答。}把目标位置切成区间两端0、\(M\)，原来的起点距离一端为\(d\)或\(M-d\)。令\(g_i\)为区间退出时间，则\(g_0=g_M=0\)，且\(g_i=1+(g_{i-1}+g_{i+1})/2\)。直接展开可检查\(i(M-i)\)满足此方程及边界。有限区间的解唯一，因此

\[\boxed{\mathbb ET=d(M-d).}\]

十二格钟面到正对面需平均\(6\cdot6=36\)步；这不是正返回均值12步。

\section{习题1.62：棋王的平稳分布与角点返回}\label{ux4e60ux98981.62ux68cbux738bux7684ux5e73ux7a33ux5206ux5e03ux4e0eux89d2ux70b9ux8fd4ux56de}

\textbf{书内依据：第1.5节无向图随机游走，第1.7节正返回均值。}棋王在\(8\times8\)棋盘上，从当前格子所有合法的相邻格中等概率选一个；不允许走出棋盘，也不包含原地不动。

\subsection{（a）逐类计数得到平稳分布}\label{aux9010ux7c7bux8ba1ux6570ux5f97ux5230ux5e73ux7a33ux5206ux5e03}

一个格子的邻居个数称为它的度\(d(x)\)。角点有3个邻居；边界上不是角点的格子有5个；内部格子有8个。三类格子数分别为4、24、36。因此度数总和是

\[D=4\cdot3+24\cdot5+36\cdot8=420.\]

对相邻格子\(x,y\)，取\(\pi_x=d(x)/D\)，则

\[\pi_xP(x,y)=\frac{d(x)}D\frac1{d(x)}=\frac1D=\pi_yP(y,x).\]

不相邻时两边都是0。故细致平衡成立，归一化也已经由\(D\)保证。答案为

\[
\boxed{\pi_x=
\begin{cases}
3/420,&x\text{为角点},\\
5/420,&x\text{为非角边界格},\\
8/420,&x\text{为内部格}.
\end{cases}}
\]

这是\textbf{每一个格子}的概率，不是每一类格子的总概率。求一类总概率还要乘该类的格子数。

\subsection{（b）角点正返回}\label{bux89d2ux70b9ux6b63ux8fd4ux56de}

有限棋盘连通，从任意角点出发的平均正返回时间为

\[\boxed{\mathbb E_xT_x=420/3=140\text{步}.}\]

均匀分布\(1/64\)不是本题答案，因为边界格与内部格的出度不同。

\subsection{变式1.62A：长方形棋盘}\label{ux53d8ux5f0f1.62aux957fux65b9ux5f62ux68cbux76d8}

\textbf{题目。}改成\(m\times n\)，其中\(m,n\ge2\)，求角点返回均值。

\textbf{解答。}四角仍度3；非角边界格有\(2m+2n-8\)个、度5；内部格有\((m-2)(n-2)\)个、度8。这些公式在某一边长为2时也成立。故

\[D=12+5(2m+2n-8)+8(m-2)(n-2)=8mn-6m-6n+4.\]

利用\(\pi_{\rm corner}=3/D\)，得到

\[\boxed{\mathbb E T_{\rm corner}=\frac{8mn-6m-6n+4}{3}.}\]

\subsection{变式1.62B：一半时间原地不动}\label{ux53d8ux5f0f1.62bux4e00ux534aux65f6ux95f4ux539fux5730ux4e0dux52a8}

\textbf{题目。}每步先以\(1/2\)原地不动，否则作原来的棋王移动。角点正返回均值是否加倍？

\textbf{解答。}新矩阵\(P'=(I+P)/2\)，原\(\pi\)满足\(\pi P'=\pi\)，所以正返回均值仍为140。原因是第一步原地不动也算一次正返回，不能忽略。

也可直接计算：原链在离开角点后返回，平均还需\(140-1=139\)个实际移动。新链每个实际移动平均用2个时钟步。条件于第一步已经移动，新正返回均值为\(1+2\cdot139=279\)。无条件均值为

\[\tfrac12\cdot1+\tfrac12\cdot279=140.\]

\subsection{变式1.62C：给斜向移动更大权重}\label{ux53d8ux5f0f1.62cux7ed9ux659cux5411ux79fbux52a8ux66f4ux5927ux6743ux91cd}

\textbf{题目。}直向相邻边权1，斜向相邻边权2，按边权比例选择，求\(8\times8\)棋盘的角点返回均值。

\textbf{解答。}每条无向边在度数总和中算两次。直向有向边数为224，斜向有向边数为196。加权总度为\(224+2\cdot196=616\)。角点有两条直向边和一条斜向边，加权度为\(1+1+2=4\)。对称边权仍保证细致平衡，故

\[\boxed{\pi_{\rm corner}=4/616,\qquad\mathbb ET_{\rm corner}=154.}\]

\section{习题1.63：皇后的平稳分布与返回时间}\label{ux4e60ux98981.63ux7687ux540eux7684ux5e73ux7a33ux5206ux5e03ux4e0eux8fd4ux56deux65f6ux95f4}

\textbf{书内依据：第1.5节图上的细致平衡，第1.7节。}每步从皇后所在格的所有合法目的格中等概率选一个。皇后可以沿行、列或斜线走任意正距离。

\subsection{（a）不能把棋王的邻居数套进来}\label{aux4e0dux80fdux628aux68cbux738bux7684ux90bbux5c45ux6570ux5957ux8fdbux6765}

行和列贡献\(7+7=14\)个目的格。对于行列编号\((i,j)\)，四条斜射线的长度分别是

\[\min(i-1,j-1),\ \min(i-1,8-j),\ \min(8-i,j-1),\ \min(8-i,8-j).\]

将其相加再加14，就是\(d(i,j)\)。设\(r\)为格子到最近边界的距离，\(r=0,1,2,3\)，整理得到\(d=21+2r\)。四层格子数为28、20、12、4，因此

\[D=28\cdot21+20\cdot23+12\cdot25+4\cdot27=1456.\]

移动关系是对称的，故与棋王题相同，\(\pi(i,j)=d(i,j)/D\)。答案为

\[\boxed{\pi(i,j)=\frac{21+2r(i,j)}{1456}.}\]

\subsection{（b）角点正返回均值}\label{bux89d2ux70b9ux6b63ux8fd4ux56deux5747ux503c}

角点可以沿同一行、同一列、一条斜线各走7个格子，所以度数21。故

\[\boxed{\mathbb E T_{\rm corner}=\frac{1456}{21}=\frac{208}{3}.}\]

皇后能走得远，不意味着角点返回均值就是棋盘格数64；关键仍是平稳概率。

\subsection{变式1.63A：中心格与整圈边界}\label{ux53d8ux5f0f1.63aux4e2dux5fc3ux683cux4e0eux6574ux5708ux8fb9ux754c}

\textbf{题目。}求一个最中心格的返回均值，以及平稳时处于最外圈的总概率。

\textbf{解答。}中心四格的度均为27，故单个中心格返回均值\(1456/27\)。最外圈共有28格，每格概率\(21/1456\)，所以总概率

\[\boxed{\frac{28\cdot21}{1456}=\frac{21}{52}.}\]

一个格子与一圈格子的概率不能相互替代。

\subsection{变式1.63B：改成车}\label{ux53d8ux5f0f1.63bux6539ux6210ux8f66}

\textbf{题目。}车每步均匀选同一行或同一列的其他格子，求平稳分布与返回均值。

\textbf{解答。}每格都恰有14个目的格，移动关系对称。因此\(\pi_x=14/(64\cdot14)=1/64\)。连通性保证唯一，返回任意格的平均时间为64。这一次均匀分布成立，是因为所有度数相等。

\subsection{\texorpdfstring{变式1.63C：任意\(N\times N\)棋盘}{变式1.63C：任意N\textbackslash times N棋盘}}\label{ux53d8ux5f0f1.63cux4efbux610fntimes-nux68cbux76d8}

\textbf{题目。}\(N\ge2\)时，求皇后从角点出发的正返回均值。

\textbf{解答。}行、列方向的有向边总数为\(2N^2(N-1)\)。一族平行斜线的长度为\(1,2,\ldots,N-1,N,N-1,\ldots,1\)。长度\(l\)的斜线贡献\(l(l-1)\)条有向边，因此一族贡献

\[2\sum_{l=1}^{N-1}l(l-1)+N(N-1)=\frac{N(N-1)(2N-1)}3.\]

两族斜线加上横竖方向，总度

\[D=\frac{2N(N-1)(5N-1)}3.\]

角点度数\(3(N-1)\)，所以

\[\boxed{\mathbb E T_{\rm corner}=\frac{2N(5N-1)}9.}\]

代入\(N=8\)，回到\(208/3\)。

\section{习题1.64：有突变的Wright--Fisher模型}\label{ux4e60ux98981.64ux6709ux7a81ux53d8ux7684wrightfisherux6a21ux578b}

\textbf{书内依据：第1.1节Wright--Fisher例（正文例1.9），第1.4节平稳分布，第1.6节有限不可约非周期链收敛。}题干的例号交叉引用不够稳定，这里按实际模型定位，不把错误例号照抄成书内依据。

设总基因数为\(N\)，当前类型\(A\)的数目为\(x\)。\(A\)以概率\(u\)突变为\(a\)，\(a\)以概率\(v\)突变为\(A\)。下一代的每个基因独立采样，成为\(A\)的概率为

\[\rho_x=(1-u)\frac{x}{N}+v\left(1-\frac{x}{N}\right)
=v+(1-u-v)\frac{x}{N}.\]

所以\(X_{n+1}\mid X_n=x\)服从参数\((N,\rho_x)\)的二项分布。

\subsection{（a）唯一平衡与一个必须说明的端点}\label{aux552fux4e00ux5e73ux8861ux4e0eux4e00ux4e2aux5fc5ux987bux8bf4ux660eux7684ux7aefux70b9}

在通常的非退化假设\(0<u,v<1\)下，\(\rho_x\)介于\(v\)与\(1-u\)之间，严格处于\((0,1)\)。于是对所有\(x,y\in\{0,\ldots,N\}\)，

\[P(x,y)=\binom Ny\rho_x^y(1-\rho_x)^{N-y}>0.\]

所有状态一步相通，且每个状态都有自环。所以链有限、不可约、非周期，有唯一平稳分布，并且从任意初始状态收敛到它。

\textbf{边界校正。}单写\(u,v>0\)还包含\(u=v=1\)（二者作为概率不超过1）。此时\(\rho_0=1\)、\(\rho_N=0\)，所以\(0\to N\to0\)确定交替。从0开始的分布不收敛。因而原题的收敛断言必须排除这一端点，不能声称仅凭严格大于0已排除所有问题。

事实上，\(0<u,v\le1\)且不同时等于1仍可成立：例如\(u=1,v<1\)时，状态0下一步的二项分布对所有状态都正，每个状态又能以正概率到0，且0有自环；\(v=1,u<1\)对称处理。

\subsection{（b）无需完整分布也能求均值}\label{bux65e0ux9700ux5b8cux6574ux5206ux5e03ux4e5fux80fdux6c42ux5747ux503c}

二项分布的均值是试验次数乘成功概率，因此

\[\mathbb E[X_{n+1}\mid X_n=x]=Nv+(1-u-v)x.\]

令\(\mu_n=\mathbb EX_n\)，对上式再取平均：

\[\mu_{n+1}=Nv+(1-u-v)\mu_n.\]

平稳均值\(\mu\)满足\(\mu=Nv+(1-u-v)\mu\)，故

\[\boxed{\sum_{x=0}^Nx\pi_x=\frac{Nv}{u+v}.}\]

从确定\(x\)开始还得到完整递推解

\[\boxed{\mu_n=\frac{Nv}{u+v}+(1-u-v)^n\left(x-\frac{Nv}{u+v}\right).}\]

在上述收敛条件下\(|1-u-v|<1\)，所以均值趋于该值。不能把``均值可求''误写成``整个平稳分布就是二项分布''。

\subsection{变式1.64A：什么时候平稳分布直接是二项分布}\label{ux53d8ux5f0f1.64aux4ec0ux4e48ux65f6ux5019ux5e73ux7a33ux5206ux5e03ux76f4ux63a5ux662fux4e8cux9879ux5206ux5e03}

\textbf{题目。}假设\(u+v=1\)，求平稳分布。

\textbf{解答。}此时\(\rho_x=v\)完全不依赖\(x\)。因此矩阵的每一行都等于\(\binom Ny v^y(1-v)^{N-y}\)，无论初始分布如何，一步后都成为这一分布。故

\[\boxed{\pi_y=\binom Ny v^y(1-v)^{N-y}.}\]

这说明二项形式在这个特殊参数关系下成立，而不是一般情况下自动成立。

\subsection{变式1.64B：平稳方差}\label{ux53d8ux5f0f1.64bux5e73ux7a33ux65b9ux5dee}

\textbf{题目。}在\(0<u,v<1\)下求平稳方差，不要求先解全部\(N+1\)个平稳概率。

\textbf{解答。}记\(a=1-u-v\)，\(p_*=v/(u+v)\)，平稳均值为\(Np_*\)。条件成功率\(\rho=v+aX/N\)满足\(\mathbb E\rho=p_*\)、\(\operatorname{Var}(\rho)=a^2\sigma^2/N^2\)。按条件均值与条件方差分解，

\[
\begin{aligned}
\sigma^2
&=\mathbb E[N\rho(1-\rho)]+\operatorname{Var}(N\rho)\\
&=Np_*(1-p_*)-\frac{a^2\sigma^2}{N}+a^2\sigma^2.
\end{aligned}
\]

移项即

\[\boxed{\sigma^2=\frac{Np_*(1-p_*)}{1-(1-1/N)a^2}.}\]

分母严格为正；当\(a=0\)时回到二项方差，另作检查。

\subsection{变式1.64C：完全没有突变}\label{ux53d8ux5f0f1.64cux5b8cux5168ux6ca1ux6709ux7a81ux53d8}

\textbf{题目。}令\(u=v=0\)，求最终全为\(A\)的概率，并说明平稳分布是否唯一。

\textbf{解答。}此时0和\(N\)都吸收。内部每个状态都有正概率下一步直接进入0或\(N\)；内部状态有限，故最终吸收概率为1。设\(h_x\)为吸收到\(N\)的概率，边界是\(h_0=0,h_N=1\)。函数\(h_x=x/N\)满足第一步方程，因为\(\mathbb E[X_{n+1}\mid X_n=x]=x\)。由吸收问题的唯一性，

\[\boxed{\mathbb P_x(\text{最终全为}A)=x/N.}\]

两个吸收状态的任意混合\(\theta\delta_N+(1-\theta)\delta_0\)都是平稳分布；内部状态是暂留的，不能有平稳质量。因此这就是全部平稳分布，不唯一。

\section{习题1.65：Ehrenfest链的均值}\label{ux4e60ux98981.65ehrenfestux94feux7684ux5747ux503c}

\textbf{书内依据：第1.1节抽球建模，条件期望，第1.2节一步递推。}共有\(N\ge1\)个球，每步均匀选一个球移到另一个盒子，\(X_n\)为左盒球数。

\subsection{（a）条件均值}\label{aux6761ux4ef6ux5747ux503c}

当前有\(x\)个球时，以\((N-x)/N\)概率增加1，以\(x/N\)概率减少1，所以

\[
\begin{aligned}
\mathbb E[X_{n+1}\mid X_n=x]
&=(x+1)\frac{N-x}{N}+(x-1)\frac{x}{N}\\
&=1+\left(1-\frac2N\right)x.
\end{aligned}
\]

端点也被包括：在0必到1，在\(N\)必到\(N-1\)。

\subsection{\texorpdfstring{（b）从\(x\)开始经过\(n\)步}{（b）从x开始经过n步}}\label{bux4ecexux5f00ux59cbux7ecfux8fc7nux6b65}

写\(\mu_n=\mathbb E_xX_n\)，则\(\mu_{n+1}=1+(1-2/N)\mu_n\)。固定值为\(N/2\)，减去它后是等比递推：

\[\mu_{n+1}-N/2=(1-2/N)(\mu_n-N/2).\]

故

\[\boxed{\mathbb E_xX_n=\frac N2+\left(1-\frac2N\right)^n\left(x-\frac N2\right).}\]

\(N\ge2\)时均值收敛到\(N/2\)；\(N=1\)时只有一个球，数目确定交替。即使\(N\ge2\)均值收敛，链的状态奇偶性仍每步改变，所以从确定初值出发，整体分布并不逐步收敛。均值收敛比概率分布收敛弱得多。

\subsection{变式1.65A：平稳分布}\label{ux53d8ux5f0f1.65aux5e73ux7a33ux5206ux5e03}

\textbf{题目。}求此链的平稳分布。

\textbf{解答。}尝试\(\pi_i=2^{-N}\binom Ni\)。相邻双向流量为

\[\pi_i\frac{N-i}{N}=\frac{\binom Ni(N-i)}{2^NN}
=\frac{\binom N{i+1}(i+1)}{2^NN}=\pi_{i+1}\frac{i+1}{N}.\]

用到的组合恒等式由阶乘定义直接约分得到。二项式定理保证\(\sum_i\pi_i=1\)，故答案为\(\operatorname{Bin}(N,1/2)\)。其存在不排除链有周期2。

\subsection{变式1.65B：均值以外的波动}\label{ux53d8ux5f0f1.65bux5747ux503cux4ee5ux5916ux7684ux6ce2ux52a8}

\textbf{题目。}求从确定\(x\)出发的方差。

\textbf{解答。}令\(Y_n=X_n-N/2\)。每步\(Y\)增减1，且条件增量均值为\(-2Y_n/N\)。因此

\[\mathbb E[Y_{n+1}^2\mid Y_n]=(1-4/N)Y_n^2+1.\]

如同原题解一阶递推，

\[\mathbb EY_n^2=\frac N4+(1-4/N)^n\left((x-N/2)^2-N/4\right).\]

再减去\((\mathbb EY_n)^2\)，得到

\[
\boxed{\operatorname{Var}_x(X_n)=\frac N4+(1-4/N)^n\left((x-N/2)^2-N/4\right)
-(1-2/N)^{2n}(x-N/2)^2.}
\]

\(n=0\)代入为0；\(N=2\)时第二项可能交替，这符合状态奇偶限制，不能强行说所有二阶量都单调收敛。

\subsection{变式1.65C：允许不移球}\label{ux53d8ux5f0f1.65cux5141ux8bb8ux4e0dux79fbux7403}

\textbf{题目。}每步以概率\(\varepsilon\in(0,1)\)执行一次原操作，否则不动，求均值与长期分布。

\textbf{解答。}条件均值成为\(\varepsilon+(1-2\varepsilon/N)x\)，故

\[\mathbb E_xX_n=N/2+(1-2\varepsilon/N)^n(x-N/2).\]

新矩阵\((1-\varepsilon)I+\varepsilon P\)保留原平稳分布。每个状态有自环且链仍连通，所以现在从任何初始状态都收敛到\(\operatorname{Bin}(N,1/2)\)。这里加自环改变的是逐步收敛，不是平稳分布本身。

\section{习题1.66：兄妹交配的六状态链}\label{ux4e60ux98981.66ux5144ux59b9ux4ea4ux914dux7684ux516dux72b6ux6001ux94fe}

\textbf{书内依据：第1.1节状态选择，第1.9节吸收概率，第1.10节退出时间。}用2、1、0分别表示\(AA,Aa,aa\)。忽略两个个体的性别后，状态顺序固定为\(22,21,20,11,10,00\)。每个后代从双亲各独立抽一个等位基因；不同后代的抽样也独立。

\subsection{（a）先算一个后代，再算两个后代}\label{aux5148ux7b97ux4e00ux4e2aux540eux4ee3ux518dux7b97ux4e24ux4e2aux540eux4ee3}

双亲是21时，后代为2、1的概率各\(1/2\)，不可能为0。因此两个后代的无序类型为22、21、11的概率分别为\(1/4,1/2,1/4\)。中间的\(1/2\)来自两种顺序\(2,1\)和\(1,2\)，不能漏掉一个。

双亲为20时，每个后代都为1，所以必到11。双亲为11时，一个后代的分布为\((1/4,1/2,1/4)\)；两个后代独立，故六种无序状态概率依次为

\[1/16,\quad2(1/4)(1/2)=1/4,\quad2(1/4)(1/4)=1/8,\quad1/4,\quad1/4,\quad1/16.\]

状态10与21对称，22和00吸收。于是完整矩阵为

\[
\boxed{P=\begin{pmatrix}
1&0&0&0&0&0\\
1/4&1/2&0&1/4&0&0\\
0&0&0&1&0&0\\
1/16&1/4&1/8&1/4&1/4&1/16\\
0&0&0&1/4&1/2&1/4\\
0&0&0&0&0&1
\end{pmatrix}.}
\]

\subsection{\texorpdfstring{（b）为什么固定为\(AA\)的概率就是当前\(A\)的比例}{（b）为什么固定为AA的概率就是当前A的比例}}\label{bux4e3aux4ec0ux4e48ux56faux5b9aux4e3aaaux7684ux6982ux7387ux5c31ux662fux5f53ux524daux7684ux6bd4ux4f8b}

每个状态包含两个二倍体个体，共四个等位基因。六状态中\(A\)的比例分别为

\[f=(1,3/4,1/2,1/2,1/4,0).\]

随机抽取不会改变\(A\)比例的条件均值，所以\(Pf=f\)；也可将上面的六行分别乘此列向量逐项检查。边界\(f(22)=1,f(00)=0\)，因此它满足固定为22的概率方程。

为什么这个方程的解就是吸收概率？从21或10下一步有\(1/4\)概率被吸收；从11下一步吸收概率\(1/8\)；从20下一步到11，再下一步有\(1/8\)概率被吸收。因此从任意暂留状态出发，接下来两步吸收概率至少\(1/8\)。连续\(2k\)步尚未吸收的概率至多\((7/8)^k\)，趋于0。有限吸收问题的边界解唯一，故

\[\boxed{\mathbb P_x(\text{最终吸收于22})=f(x).}\]

这里用的是``下一代平均比例不变''，不是``每条样本路径的比例恒定''。样本路径上的比例会变化。

\subsection{（c）所有起点的平均吸收时间}\label{cux6240ux6709ux8d77ux70b9ux7684ux5e73ux5747ux5438ux6536ux65f6ux95f4}

令\(g(x)=\mathbb E_xT\)，到达22或00即停止，所以这两个值为0。对称性给出\(g(21)=g(10)=a\)，另记\(g(20)=b,g(11)=c\)。第一步方程是

\[
\begin{aligned}
a&=1+\tfrac12a+\tfrac14c,\\
b&=1+c,\\
c&=1+\tfrac14a+\tfrac18b+\tfrac14c+\tfrac14a.
\end{aligned}
\]

第一式给\(a=2+c/2\)，第二式给\(b=1+c\)。代入第三式，

\[c=1+\tfrac12(2+c/2)+\tfrac18(1+c)+\tfrac14c
=\tfrac{17}{8}+\tfrac58c.\]

所以\(c=17/3\)，随后\(a=29/6,b=20/3\)。全部答案按原状态顺序为

\[\boxed{g=(0,29/6,20/3,17/3,29/6,0).}\]

状态20当前看起来没有杂合子，但下一代必然全为杂合子；所以它的平均吸收时间反而比11多一步。这正好检查\(b=1+c\)。

\subsection{变式1.66A：为什么可以忽略性别}\label{ux53d8ux5f0f1.66aux4e3aux4ec0ux4e48ux53efux4ee5ux5ffdux7565ux6027ux522b}

\textbf{题目。}原先保留父、母顺序有九状态，证明合并成六状态后仍然是马尔可夫链。

\textbf{解答。}给定父母类型\(a,b\)，设一个后代的三类型概率是\(q_2,q_1,q_0\)。交换父母顺序不改变\(q\)。下一对后代的有序类型\((c,d)\)概率为\(q_cq_d\)；合并后，相同类型\(cc\)概率\(q_c^2\)，不同类型\(cd\)概率\(2q_cq_d\)。这些概率仅依赖父母的无序类型，不依赖被丢弃的性别顺序，更不依赖历史。因此同一合并类内的任何状态都有相同的类间转移概率，这恰好保证六状态过程仍是马尔可夫链。

\subsection{\texorpdfstring{变式1.66B：条件于最终固定为\(AA\)，还要等多久}{变式1.66B：条件于最终固定为AA，还要等多久}}\label{ux53d8ux5f0f1.66bux6761ux4ef6ux4e8eux6700ux7ec8ux56faux5b9aux4e3aaaux8fd8ux8981ux7b49ux591aux4e45}

\textbf{题目。}从21开始，已知最终吸收到22，求平均吸收时间。

\textbf{解答。}令\(u_x=\mathbb E_x[T\mathbf1_{\{\text{吸收22}\}}]\)，不能直接将无条件均值除以成功概率。第一步已经用掉的一代仅在最终成功的路径上计入，所以暂留状态上\(u=h+Qu\)，其中\(h=(3/4,1/2,1/2,1/4)^{\mathsf T}\)，状态顺序\(21,20,11,10\)。

具体方程为

\[
\begin{aligned}
u_{21}&=3/4+u_{21}/2+u_{11}/4,\\
u_{20}&=1/2+u_{11},\\
u_{11}&=1/2+u_{21}/4+u_{20}/8+u_{11}/4+u_{10}/4,\\
u_{10}&=1/4+u_{11}/4+u_{10}/2.
\end{aligned}
\]

首、末两式给\(u_{21}=3/2+u_{11}/2\)，\(u_{10}=1/2+u_{11}/2\)。代入中间方程得\(u_{11}=17/6\)，因此\(u_{21}=35/12\)。最后才除以\(h_{21}=3/4\)：

\[\boxed{\mathbb E_{21}[T\mid\text{吸收22}]=\frac{35/12}{3/4}=\frac{35}{9}.}\]

\subsection{变式1.66C：吸收前有多少代处于11}\label{ux53d8ux5f0f1.66cux5438ux6536ux524dux6709ux591aux5c11ux4ee3ux5904ux4e8e11}

\textbf{题目。}从20开始，求吸收之前处于11的总次数的期望。起始时刻若是11也应计一次。

\textbf{解答。}记期望为\(v_x\)，在11每次立即获得奖励1，其他暂留状态奖励0。因此\(v_{21}=v_{10}=v_{11}/2\)，\(v_{20}=v_{11}\)，且

\[v_{11}=1+\tfrac14v_{21}+\tfrac18v_{20}+\tfrac14v_{11}+\tfrac14v_{10}
=1+\tfrac58v_{11}.\]

所以\(v_{11}=8/3\)，原问题从20出发的答案也为\(\boxed{8/3}\)。这不是总吸收时间：其他类型的世代不计入这个奖励。

\section{习题1.67：六面骰集齐全部点数}\label{ux4e60ux98981.67ux516dux9762ux9ab0ux96c6ux9f50ux5168ux90e8ux70b9ux6570}

\textbf{书内依据：第1.1节状态压缩，第1.10节第一步时间方程，附录几何分布的均值。}\(X_n\)记录前\(n\)次掷骰见过多少种点数，取\(X_0=0\)。

\subsection{（a）转移概率}\label{aux8f6cux79fbux6982ux7387-1}

当前见过\(k\)种时，下次出现旧点数的概率是\(k/6\)，出现新点数的概率是\((6-k)/6\)。公平性使这些概率不依赖具体见过哪些点数，所以只记录数量足够。故

\[\boxed{P(k,k)=k/6,\qquad P(k,k+1)=(6-k)/6\quad(0\le k\le5),\quad P(6,6)=1.}\]

其余概率为0。每一步不可能同时新增两种。

\subsection{（b）平均需要多少次}\label{bux5e73ux5747ux9700ux8981ux591aux5c11ux6b21}

记\(g_k\)为已有\(k\)种后，还需多少次才能集齐的均值。边界\(g_6=0\)。第一步方程

\[g_k=1+\frac{k}{6}g_k+\frac{6-k}{6}g_{k+1}\]

给出\(g_k-g_{k+1}=6/(6-k)\)。从\(k=0\)加到5，所有中间\(g_k\)抵消：

\[\boxed{\mathbb E_0T=g_0=6\left(1+\frac12+\frac13+\frac14+\frac15+\frac16\right)=\frac{147}{10}=14.7.}\]

最后一种平均就要等6次，因此``共有六种，所以平均六次''漏掉了重复点数。

\subsection{变式1.67A：在给定次数内集齐的概率}\label{ux53d8ux5f0f1.67aux5728ux7ed9ux5b9aux6b21ux6570ux5185ux96c6ux9f50ux7684ux6982ux7387}

\textbf{题目。}求\(n\ge1\)次以内集齐的概率，并算恰好前六次集齐的概率。

\textbf{解答。}设\(A_j\)表示点数\(j\)没有出现。要求的是所有\(A_j\)都不发生。固定\(j\)种点数均未出现的概率为\(((6-j)/6)^n\)。将遗漏事件用容斥相加，得到

\[\boxed{\mathbb P(T\le n)=\sum_{j=0}^6(-1)^j\binom6j\left(\frac{6-j}{6}\right)^n.}\]

\(n=6\)时必须每种恰出现一次，共\(6!\)种序列、总共\(6^6\)种等可能序列，所以

\[\boxed{\mathbb P(T\le6)=6!/6^6=5/324}.\]

这与容斥式相符。

\subsection{变式1.67B：只收集四种}\label{ux53d8ux5f0f1.67bux53eaux6536ux96c6ux56dbux79cd}

\textbf{题目。}至少见过四种点数平均需要多少次？

\textbf{解答。}停止在4而非6，分别等待第1、2、3、4种新点数，其成功概率为\(1,5/6,4/6,3/6\)。期望可直接相加，不要求这些等待时间先证明独立。因此

\[\boxed{\mathbb ET_4=1+\frac65+\frac64+\frac63=\frac{57}{10}.}\]

\subsection{变式1.67C：不等概率的三种物品}\label{ux53d8ux5f0f1.67cux4e0dux7b49ux6982ux7387ux7684ux4e09ux79cdux7269ux54c1}

\textbf{题目。}每次独立得到三种物品之一，概率为\(1/2,1/3,1/6\)，求集齐均值。

\textbf{解答。}只记录``已见种数''不再足够，因为缺少罕见种与缺少常见种的等待时间不同。直接对缺失事件使用容斥：对\(n\ge0\)，未集齐概率为\(\sum_i(1-p_i)^n-\sum_{i<j}(1-p_i-p_j)^n+(1-\sum_i p_i)^n\)。在\(n=0\)处各事件的概率都取1；即这里的\(0^0\)按事件概率解释为1。

利用\(\mathbb ET=\sum_{n\ge0}\mathbb P(T>n)\)和等比级数，得到

\[\mathbb ET=\sum_i\frac1{p_i}-\sum_{i<j}\frac1{p_i+p_j}+\frac1{p_1+p_2+p_3}.\]

代入三概率：\(2+3+6-(6/5+3/2+2)+1=\boxed{73/10}\)。

\section{\texorpdfstring{习题1.68：集齐\(N\)种物品的均值、方差与渐近量级}{习题1.68：集齐N种物品的均值、方差与渐近量级}}\label{ux4e60ux98981.68ux96c6ux9f50nux79cdux7269ux54c1ux7684ux5747ux503cux65b9ux5deeux4e0eux6e10ux8fd1ux91cfux7ea7}

\textbf{书内依据：第1.10节等待时间分解，附录独立变量方差与几何分布，微积分中的积分比较。}每次独立、等概率取得\(N\)种物品之一，\(T_k\)为首次集齐\(k\)种的时刻，令\(T_0=0\)。

\subsection{第一步：把总时间拆成几何等待}\label{ux7b2cux4e00ux6b65ux628aux603bux65f6ux95f4ux62c6ux6210ux51e0ux4f55ux7b49ux5f85}

设\(G_k=T_{k+1}-T_k\)，\(0\le k\le N-1\)。已收集\(k\)种后，每次新增一种的概率为\((N-k)/N\)。因此\(G_k\)取值\(1,2,\ldots\)，为成功概率\(a_k=(N-k)/N\)的几何变量。

还需解释方差计算使用的独立性：给定收集到\(k\)种以前的全部历史，后面的抽取仍是新的独立试验；而\(G_k\)的几何参数只由\(k\)决定，不依赖此前收到了哪些种、花了多久。所以\(G_k\)的条件分布不随此前等待时间变化。逐项使用条件概率乘法，便得这些\(G_k\)相互独立。

\subsection{第二步：精确均值和方差}\label{ux7b2cux4e8cux6b65ux7cbeux786eux5747ux503cux548cux65b9ux5dee}

几何变量的均值是\(1/a\)，方差是\((1-a)/a^2\)，且\(T_N=\sum_{k=0}^{N-1}G_k\)。令\(H_N=\sum_{j=1}^N1/j\)，\(H_N^{(2)}=\sum_{j=1}^N1/j^2\)。则

\[\boxed{\mathbb ET_N=NH_N,\qquad\operatorname{Var}(T_N)=N^2H_N^{(2)}-NH_N.}\]

方差第二项是减去\(NH_N\)，不能漏掉。它来自每个几何方差中的\(-1/a_k\)。

\subsection{第三步：不用模糊的``约等于''代替极限}\label{ux7b2cux4e09ux6b65ux4e0dux7528ux6a21ux7ccaux7684ux7ea6ux7b49ux4e8eux4ee3ux66ffux6781ux9650}

函数\(1/x\)递减，矩形面积与积分比较给

\[\log(N+1)\le H_N\le1+\log N.\]

两边除以\(\log N\)，得到\(H_N/\log N\to1\)，故

\[\boxed{\frac{\mathbb ET_N}{N\log N}\longrightarrow1.}\]

再看方差：

\[\frac{\operatorname{Var}(T_N)}{N^2}=H_N^{(2)}-\frac{H_N}{N}.\]

第二项趋于0；第一项递增到有限的正数\(\sum_{j=1}^\infty1/j^2\)，因为尾部可由\(\int_N^\infty x^{-2}\,dx=1/N\)控制。因此

\[\boxed{\frac{\operatorname{Var}(T_N)}{N^2}\longrightarrow\sum_{j=1}^\infty\frac1{j^2}.}\]

这就是原题的两个渐近结论，不需要先计算这个级数的闭式。

\subsection{变式1.68A：95\%把握收齐需要多少次}\label{ux53d8ux5f0f1.68a95ux628aux63e1ux6536ux9f50ux9700ux8981ux591aux5c11ux6b21}

\textbf{题目。}给出一个易计算、对任意\(N\)都成立的充分次数。

\textbf{解答。}固定一种还没出现的概率为\((1-1/N)^n\)。至少有一种未出现的概率不超过这些事件概率之和，故

\[\mathbb P(T_N>n)\le N(1-1/N)^n\le Ne^{-n/N}.\]

最后一步用\(1-x\le e^{-x}\)，可由函数\(e^{-x}-(1-x)\)的导数证明。只要

\[\boxed{n\ge\left\lceil N(\log N+\log20)\right\rceil,}\]

失败概率就不超过\(1/20\)。这是保证，不声称是最小次数。

\subsection{变式1.68B：收集固定比例而非全部}\label{ux53d8ux5f0f1.68bux6536ux96c6ux56faux5b9aux6bd4ux4f8bux800cux975eux5168ux90e8}

\textbf{题目。}求首次收齐\(k\le N\)种的期望；若\(k/N\to\theta\in(0,1)\)，求其渐近量级。

\textbf{解答。}只需相加前\(k\)个等待阶段，

\[\boxed{\mathbb ET_k=N(H_N-H_{N-k}).}\]

当\(k/N\to\theta<1\)时，积分比较给\(H_N-H_{N-k}=\log(N/(N-k))+o(1)\)。因此

\[\boxed{\mathbb ET_k/N\longrightarrow-\log(1-\theta).}\]

收固定比例只需要\(N\)量级；集齐最后少数稀缺种才把总期望推高到\(N\log N\)量级。

\subsection{变式1.68C：随机总时间是否接近其均值}\label{ux53d8ux5f0f1.68cux968fux673aux603bux65f6ux95f4ux662fux5426ux63a5ux8fd1ux5176ux5747ux503c}

\textbf{题目。}证明\(T_N/(N\log N)\)依概率趋于1。

\textbf{解答。}由\(\sum_{j\ge1}j^{-2}\le1+\int_1^\infty x^{-2}dx=2\)，有\(\operatorname{Var}(T_N)\le2N^2\)。对非负变量\((T_N-\mathbb ET_N)^2\)使用期望上界，得到

\[\mathbb P\left(\left|\frac{T_N}{\mathbb ET_N}-1\right|>\varepsilon\right)
\le\frac{2}{\varepsilon^2H_N^2}\longrightarrow0.\]

又\(\mathbb ET_N/(N\log N)\to1\)，两步合并便得结论。这里不是说每条样本都等于期望，而是相对偏差超过固定阈值的概率趋于0。

\section{习题1.69：严格下降链与独立访问指标}\label{ux4e60ux98981.69ux4e25ux683cux4e0bux964dux94feux4e0eux72ecux7acbux8bbfux95eeux6307ux6807}

\textbf{书内依据：第1.10节第一步时间方程，条件概率与独立性。}状态\(1,\ldots,n\)，1吸收；从\(i>1\)均匀跳到\(1,\ldots,i-1\)之一。每个目的地概率应为\(1/(i-1)\)；题面分式排版若看成\(1/i-1\)会得到负数，不是本模型。

本题是题设的随机下降模型。以下结果不自动等于任意实际单纯形算法的运行时间定理。

\subsection{（a）相邻两个期望相减}\label{aux76f8ux90bbux4e24ux4e2aux671fux671bux76f8ux51cf}

设\(e_i\)为从\(i\)到1的平均步数，\(e_1=0\)。从\(i>1\)出发，

\[e_i=1+\frac1{i-1}\sum_{j=1}^{i-1}e_j.\]

\(i=2\)时\(e_2=1\)。对\(i\ge3\)，分别乘分母再相减：

\[
\begin{aligned}
(i-1)e_i&=(i-1)+\sum_{j=1}^{i-1}e_j,\\
(i-2)e_{i-1}&=(i-2)+\sum_{j=1}^{i-2}e_j.
\end{aligned}
\]

相减得\((i-1)e_i-(i-2)e_{i-1}=1+e_{i-1}\)，所以

\[e_i-e_{i-1}=\frac1{i-1}.\]

从2加到\(i\)，得到

\[\boxed{e_i=1+\frac12+\cdots+\frac1{i-1}=H_{i-1}.}\]

\subsection{（b）为什么访问不同状态的指标独立}\label{bux4e3aux4ec0ux4e48ux8bbfux95eeux4e0dux540cux72b6ux6001ux7684ux6307ux6807ux72ecux7acb}

从\(n\)出发，令\(I_j\)表示是否访问\(j\)，\(I_n=1\)。固定\(j<n\)，并给定所有高于\(j\)的状态是否被访问。因为路径严格下降，这等于给定了所有被访问的高状态及其顺序。

设其中最后一个高于\(j\)的已访问状态为\(k\)。给定这段高状态路径后，从\(k\)离开时下一跳已经知道不能落在\(j+1,\ldots,k-1\)，否则这些高状态也会被记录为访问。原来从\(k\)均匀选\(1,\ldots,k-1\)；在条件``选择不超过\(j\)''下，它仍均匀选\(1,\ldots,j\)。因此

\[\boxed{\mathbb P(I_j=1\mid I_{j+1},\ldots,I_n)=\frac1j.}\]

这个等式只对正概率的给定高状态组合使用。若第一次跳到小于\(j\)的位置，之后严格下降，不可能补访\(j\)；若跳到\(j\)，则访问成功。故条件概率确实恰好是上面的一个选中概率。

对\(I_{n-1},I_{n-2},\ldots,I_1\)依次使用乘法公式，每一步成功概率\(1/j\)都不随已给定的高指标变化，失败概率也固定为\(1-1/j\)。于是任意0、1组合的联合概率等于各自边际概率之积，这证明了\(I_1,\ldots,I_{n-1}\)相互独立。\(I_1\equiv1\)是允许的退化独立变量。

下降到1的步数等于访问的较低状态数：\(T_1=\sum_{j=1}^{n-1}I_j\)。所以

\[\mathbb E_nT_1=\sum_{j=1}^{n-1}\mathbb EI_j=\sum_{j=1}^{n-1}\frac1j,\]

给出第二个证明。求这个期望只需线性性；独立性是额外结论，可用于下面的方差。

\subsection{变式1.69A：四状态时的完整分布}\label{ux53d8ux5f0f1.69aux56dbux72b6ux6001ux65f6ux7684ux5b8cux6574ux5206ux5e03}

\textbf{题目。}从4开始，求\(T_1\)的全部概率。

\textbf{解答。}\(T_1=1+I_2+I_3\)，后两个独立、成功概率为\(1/2,1/3\)。因此

\[\mathbb P(T_1=1)=\tfrac12\cdot\tfrac23=\tfrac13,\]

\[\mathbb P(T_1=2)=\tfrac12\cdot\tfrac23+\tfrac12\cdot\tfrac13=\tfrac12,\]

\[\mathbb P(T_1=3)=\tfrac12\cdot\tfrac13=\tfrac16.\]

期望\(1/3+1+1/2=11/6\)，等于\(H_3\)。

\subsection{变式1.69B：任意起点的方差和生成函数}\label{ux53d8ux5f0f1.69bux4efbux610fux8d77ux70b9ux7684ux65b9ux5deeux548cux751fux6210ux51fdux6570}

\textbf{题目。}求\(\operatorname{Var}_n(T_1)\)与\(\mathbb E_n[s^{T_1}]\)。

\textbf{解答。}独立伯努利变量方差相加，得到

\[\boxed{\operatorname{Var}_n(T_1)=H_{n-1}-H_{n-1}^{(2)}.}\]

又\(s^{\sum I_j}=\prod s^{I_j}\)，独立性使期望分解，所以

\[\boxed{\mathbb E_n[s^{T_1}]=\prod_{j=1}^{n-1}\left(1+\frac{s-1}{j}\right).}\]

这是一个有限多项式，展开后\(s^k\)的系数就是\(k\)步结束的概率，不需要额外收敛条件。

\subsection{变式1.69C：目标状态具有不同权重}\label{ux53d8ux5f0f1.69cux76eeux6807ux72b6ux6001ux5177ux6709ux4e0dux540cux6743ux91cd}

\textbf{题目。}给定正权重\(w_j\)，从\(i\)以概率\(w_j/W_{i-1}\)跳到\(j<i\)，其中\(W_k=\sum_{j=1}^kw_j\)。求访问指标和平均下降步数。

\textbf{解答。}重复原题的条件论证：给定高状态路径后，最后一次从高状态向不超过\(j\)的区域跳跃时，条件概率按\(w_1,\ldots,w_j\)重新归一化。因此

\[\mathbb P(I_j=1\mid I_{j+1},\ldots,I_n)=w_j/W_j.\]

条件分布与高指标无关，故指标仍独立，且

\[\boxed{\mathbb E_nT_1=\sum_{j=1}^{n-1}\frac{w_j}{W_j}.}\]

等权\(w_j=1\)回到原题，但一般权重下期望不再是调和数。

\section{习题1.70：一般生灭链的常返判别}\label{ux4e60ux98981.70ux4e00ux822cux751fux706dux94feux7684ux5e38ux8fd4ux5224ux522b}

\textbf{书内依据：第1.9节首达概率，第1.11节无穷状态链。}状态为\(0,1,2,\ldots\)。从\(x\)向上、向下、原地的概率分别为\(p_x,q_x,r_x\)，其中\(x\ge1\)时\(p_x+q_x+r_x=1\)；在0处只有向上和原地，\(p_0+r_0=1\)。

\textbf{适用条件。}下述比值公式采用通常的不可约生灭链假设：\(p_0>0\)，且\(x\ge1\)时\(p_x,q_x>0\)。如果允许\(p_0=0\)，0本身吸收、必然常返，而题目的级数可任意变化；因此不能省略这些假设后仍宣称同一``当且仅当''。

\subsection{第一步：有限区间的边界问题}\label{ux7b2cux4e00ux6b65ux6709ux9650ux533aux95f4ux7684ux8fb9ux754cux95eeux9898}

记\(V_y=\min\{n\ge0:X_n=y\}\)，并设

\[h_N(x)=\mathbb P_x(V_N<V_0),\qquad0\le x\le N.\]

边界是\(h_N(0)=0,h_N(N)=1\)。内部按第一步分类：

\[h_N(x)=p_xh_N(x+1)+r_xh_N(x)+q_xh_N(x-1).\]

将\(r_x=1-p_x-q_x\)代入并整理，

\[p_x\bigl(h_N(x+1)-h_N(x)\bigr)=q_x\bigl(h_N(x)-h_N(x-1)\bigr).\]

\subsection{第二步：把二阶方程降成一阶乘法}\label{ux7b2cux4e8cux6b65ux628aux4e8cux9636ux65b9ux7a0bux964dux6210ux4e00ux9636ux4e58ux6cd5}

定义相邻差\(\Delta_x=h_N(x)-h_N(x-1)\)，则

\[\Delta_{x+1}=\frac{q_x}{p_x}\Delta_x.\]

记\(s_0=1\)，\(s_k=\prod_{j=1}^k(q_j/p_j)\)。因\(h_N(1)=\Delta_1=c_N\)，连续相乘给\(\Delta_{k+1}=c_Ns_k\)，再相加：

\[h_N(x)=c_N\sum_{k=0}^{x-1}s_k.\]

令\(x=N\)，利用\(h_N(N)=1\)，得到

\[\boxed{c_N=\frac1{\sum_{k=0}^{N-1}s_k},\qquad
h_N(x)=\frac{\sum_{k=0}^{x-1}s_k}{\sum_{k=0}^{N-1}s_k}.}\]

这一步解释了原题中的空乘积：\(s_0=1\)对应最初的差\(\Delta_1\)，不是可以随意丢掉的一项。

\subsection{第三步：为什么可以令上边界趋于无穷}\label{ux7b2cux4e09ux6b65ux4e3aux4ec0ux4e48ux53efux4ee5ux4ee4ux4e0aux8fb9ux754cux8d8bux4e8eux65e0ux7a77}

从1开始，事件\(\{V_N<V_0\}\)随\(N\)增加而递减。若始终不返回0又永远停留在某个有限集合\(\{1,\ldots,M\}\)，则在每段长度\(M\)的时间里，存在一个统一正概率沿连续向下路径到0。这个概率是有限多个正路径概率的最小值。因此永远不去0且始终留在这个集合的概率为0。

所以``不返回0''的路径，除零概率例外，必定先后达到每个高度\(N\)。反过来，若在到0之前达到每个\(N\)，就不可能在有限时刻到0。由递减事件概率的连续性，

\[\mathbb P_1(V_0=\infty)=\lim_{N\to\infty}h_N(1)
=\frac1{\sum_{k=0}^{\infty}s_k},\]

分母为无穷时右边解释为0。从0开始，第一步原地就已正返回；只有以概率\(p_0\)走到1后再不回来，才会永不正返回。因此

\[\mathbb P_0(T_0=\infty)=\frac{p_0}{\sum_{k=0}^{\infty}s_k}.\]

由于\(p_0>0\)，终于得到

\[\boxed{0\text{常返}\iff\sum_{k=0}^{\infty}\prod_{j=1}^k\frac{q_j}{p_j}=\infty.}\]

\textbf{概念检查：}这只判定返回概率是否为1，没有判定返回时间均值是否有限；后者才区分正常返与零常返。

\subsection{变式1.70A：固定左右概率}\label{ux53d8ux5f0f1.70aux56faux5b9aux5de6ux53f3ux6982ux7387}

\textbf{题目。}内部\(p_x=p,q_x=q=1-p\)，其中\(0<p<1\)，边界\(p_0=1\)。判定常返，并在向上偏时求从\(x\)逃离0的概率。

\textbf{解答。}\(s_k=(q/p)^k\)。级数发散当且仅当\(q/p\ge1\)，故\(p\le1/2\)常返；\(p>1/2\)暂留。后一情形有限首达公式中分母趋于\(1/(1-q/p)\)，于是\(\mathbb P_x(V_0<\infty)=(q/p)^x\)，所以

\[\boxed{\mathbb P_x(V_0=\infty)=1-(q/p)^x\quad(p>1/2).}\]

向下偏\(p<1/2\)时还可求平稳权重：\(w_0=1\)，\(w_n=q^{-1}(p/q)^{n-1}\)，其和有限，因此正常返。\(p=1/2\)时权重恒定尾部不能归一化，为零常返。

\subsection{变式1.70B：减慢时钟会不会改变常返}\label{ux53d8ux5f0f1.70bux51cfux6162ux65f6ux949fux4f1aux4e0dux4f1aux6539ux53d8ux5e38ux8fd4}

\textbf{题目。}在状态\(x\)以概率\(\varepsilon_x\in(0,1]\)执行原转移，否则原地停留，讨论常返与正常返。

\textbf{解答。}新的向上、向下概率同时乘\(\varepsilon_x\)，比值\(q_x/p_x\)不变，所以首达概率与常返性不变。然而平稳权重变为\(w'_x=w_x/\varepsilon_x\)，因为

\[w'_xp'_x=w_xp_x=w_{x+1}q_{x+1}=w'_{x+1}q'_{x+1}.\]

它可能从可求和变为不可求和。例如原链内部\(p=1/3,q=2/3,p_0=1\)，有\(w_0=1,w_n=(3/2)(1/2)^{n-1}\)，总和4，正常返。改用\(\varepsilon_0=1,\varepsilon_n=2^{-n}\)，则\(n\ge1\)时\(w'_n=3\)，不能归一化。链仍常返，却变成零常返。极远处等待太久，便能使平均返回时间无穷。

\subsection{变式1.70C：加入一半停留的有限吸收区间}\label{ux53d8ux5f0f1.70cux52a0ux5165ux4e00ux534aux505cux7559ux7684ux6709ux9650ux5438ux6536ux533aux95f4}

\textbf{题目。}在\(0,\ldots,N\)上，0与\(N\)吸收，内部向上和向下各\(1/4\)、原地\(1/2\)。求先到\(N\)的概率和平均吸收时间。

\textbf{解答。}首达方程消去原地项后相邻差仍相等，所以\(h_i=i/N\)。时间方程为\(g_i=1+(g_{i-1}+g_{i+1})/4+g_i/2\)，即

\[g_{i+1}-2g_i+g_{i-1}=-4.\]

函数\(2i(N-i)\)的二阶差恰为\(-4\)，且两端为0。因此

\[\boxed{h_i=i/N,\qquad g_i=2i(N-i).}\]

停留不改变从哪一端退出，却将平均时间加倍。

\section{习题1.71：一个临界幂次决定常返的例子}\label{ux4e60ux98981.71ux4e00ux4e2aux4e34ux754cux5e42ux6b21ux51b3ux5b9aux5e38ux8fd4ux7684ux4f8bux5b50}

\textbf{书内依据：习题1.70的级数判别，第1.11节，积分比较。}采用\(c>0\)；内部

\[p_x=\frac12,\qquad q_x=\frac12e^{-cx^{-\alpha}},\qquad
r_x=\frac12-q_x,\quad x\ge1,\]

并令\(p_0=1\)。若\(c<0\)，则\(q_x>1/2,r_x<0\)，不是合法链。\(c=0\)另作边界处理。

\subsection{第一步：把乘积变成指数中的和}\label{ux7b2cux4e00ux6b65ux628aux4e58ux79efux53d8ux6210ux6307ux6570ux4e2dux7684ux548c}

因为\(q_x/p_x=e^{-cx^{-\alpha}}\)，故

\[s_k=\exp\left(-c\sum_{j=1}^kj^{-\alpha}\right).\]

根据上一题，只需判断\(\sum_{k\ge0}s_k\)是否发散。

\subsection{\texorpdfstring{第二步：\(\alpha>1\)}{第二步：\textbackslash alpha\textgreater1}}\label{ux7b2cux4e8cux6b65alpha1}

级数\(\sum_{j\ge1}j^{-\alpha}\)有限，设其和为\(C\)。则\(s_k\ge e^{-cC}>0\)。无穷多个项都有同一个正下界，所以\(\sum s_k=\infty\)。因此链常返。

\subsection{\texorpdfstring{第三步：\(\alpha=1\)}{第三步：\textbackslash alpha=1}}\label{ux7b2cux4e09ux6b65alpha1}

调和数满足\(\log(k+1)\le H_k\le1+\log k\)，故对\(k\ge1\)，

\[e^{-c}k^{-c}\le e^{-cH_k}\le(k+1)^{-c}.\]

因此\(\sum s_k\)与\(\sum k^{-c}\)具有相同的敛散性。积分\(\int_1^\infty x^{-c}dx\)在\(c\le1\)时发散，在\(c>1\)时收敛，得到常返当且仅当\(c\le1\)。

\subsection{\texorpdfstring{第四步：\(\alpha<1\)}{第四步：\textbackslash alpha\textless1}}\label{ux7b2cux56dbux6b65alpha1}

令\(\beta=1-\alpha>0\)。存在只依赖\(\alpha\)的正常数\(A\)，使大\(k\)时

\[\sum_{j=1}^kj^{-\alpha}\ge A k^\beta.\]

当\(0\le\alpha<1\)，只取\(j\)介于\(k/2\)和\(k\)的约\(k/2\)项，每项至少\(k^{-\alpha}\)；当\(\alpha<0\)，这些项每项至少\((k/2)^{-\alpha}\)，同样得到该结论。因此\(s_k\le e^{-cAk^\beta}\)。

又\(k^\beta/\log k\to\infty\)（可对实变量用洛必达法则验证），故充分大时\(cAk^\beta\ge2\log k\)，于是\(s_k\le k^{-2}\)。比较可知\(\sum s_k<\infty\)，所以暂留。

综上，在\(c>0\)条件下，

\[\boxed{\text{常返}\iff\alpha>1\ \text{或}\ (\alpha=1\text{且}c\le1).}\]

其余情形暂留。\(c=0\)时是反射公平游走，对所有\(\alpha\)都常返，不能机械套用\(c>0\)下的``其余暂留''。

\subsection{变式1.71A：常返的那些情形是否正常返}\label{ux53d8ux5f0f1.71aux5e38ux8fd4ux7684ux90a3ux4e9bux60c5ux5f62ux662fux5426ux6b63ux5e38ux8fd4}

\textbf{题目。}将原题中已经判为常返的参数，进一步分类。

\textbf{解答。}相邻细致平衡给权重\(w_0=1\)，而\(n\ge1\)时

\[w_n=\frac{p_0p_1\cdots p_{n-1}}{q_1q_2\cdots q_n}
=2\exp\left(c\sum_{j=1}^nj^{-\alpha}\right)\ge2.\]

所以\(\sum_nw_n=\infty\)，不能归一化。任何平稳分布都必须满足这种相邻流量等式：对有限集合\(\{0,\ldots,n\}\)求平衡，只有边\(n\leftrightarrow n+1\)跨越边界。因此不存在其他绕开此权重的平稳概率。由不可约链的正常返判别，所有常返参数实际上都是\textbf{零常返}。

\subsection{变式1.71B：临界零参数为何平均返回时间无穷}\label{ux53d8ux5f0f1.71bux4e34ux754cux96f6ux53c2ux6570ux4e3aux4f55ux5e73ux5747ux8fd4ux56deux65f6ux95f4ux65e0ux7a77}

\textbf{题目。}在\(c=0\)时，不借助求平稳分布，说明返回0的均值无穷。

\textbf{解答。}从1出发，先到0或\(N\)的平均时间是\(1(N-1)=N-1\)，由公平区间退出公式得到。它不超过真正到0的时间，因为我们允许在\(N\)提前停下。因此\(\mathbb E_1V_0\ge N-1\)对所有\(N\)成立，只能有\(\mathbb E_1V_0=\infty\)。从0先确定到1，再返回也有无穷均值。与此同时常返判别已说明返回概率为1。这是``几乎必然有限''与``有限期望''不同的具体例子。

\subsection{变式1.71C：再加一个对数项会怎样}\label{ux53d8ux5f0f1.71cux518dux52a0ux4e00ux4e2aux5bf9ux6570ux9879ux4f1aux600eux6837}

\textbf{题目。}设\(c,d\ge0\)，内部\(p_x=1/2\)，并令

\[\frac{q_x}{p_x}=\exp\left(-\frac{c}{x+1}-\frac{d}{(x+1)\log(x+1)}\right),\qquad p_0=1.\]

求常返条件。

\textbf{解答。}正参数保证\(q_x\le1/2\)，可用停留补足行和。两个正递减函数的和与积分之差一致有界，所以

\[\sum_{x=1}^k\frac1{x+1}=\log k+O(1),\qquad
\sum_{x=1}^k\frac1{(x+1)\log(x+1)}=\log\log k+O(1).\]

这里\(O(1)\)仅表示有一个不随\(k\)增长的界。指数化后，\(s_k\)夹在\(k^{-c}(\log k)^{-d}\)的两个正常数倍之间。若\(c<1\)，此级数发散；若\(c>1\)，收敛。可用任意小\(\eta>0\)和\((\log k)^d\le k^\eta\)作幂级数比较。\(c=1\)时用换元\(t=\log x\)，积分变成\(\int t^{-d}dt\)，发散当且仅当\(d\le1\)。因此

\[\boxed{\text{常返}\iff c<1\ \text{或}\ (c=1\text{且}d\le1).}\]

这个变式说明临界幂次为1时，低一阶的对数修正仍能改变答案。

\section{习题1.72：可归一化的无穷状态平稳分布}\label{ux4e60ux98981.72ux53efux5f52ux4e00ux5316ux7684ux65e0ux7a77ux72b6ux6001ux5e73ux7a33ux5206ux5e03}

\textbf{书内依据：第1.5节细致平衡，第1.11节正常返与平稳概率的关系。}本题

\[p_m=\frac{m+1}{2(m+2)},\quad m\ge0;\qquad
q_m=\frac{m+3}{2(m+2)},\quad m\ge1,\]

且\(P(0,0)=3/4\)。内部\(p_m+q_m=1\)。

\subsection{第一步：注意相邻索引不同}\label{ux7b2cux4e00ux6b65ux6ce8ux610fux76f8ux90bbux7d22ux5f15ux4e0dux540c}

细致平衡是\(\pi_mp_m=\pi_{m+1}q_{m+1}\)，因此

\[\frac{\pi_{m+1}}{\pi_m}=\frac{p_m}{q_{m+1}}
=\frac{(m+1)(m+3)}{(m+2)(m+4)}.\]

这里是\(q_{m+1}\)，不是\(q_m\)；与1.70中首达差分所用的\(q_m/p_m\)不要混淆。

\subsection{第二步：连乘并归一化}\label{ux7b2cux4e8cux6b65ux8fdeux4e58ux5e76ux5f52ux4e00ux5316}

取\(w_0=1\)，连乘望远镜约分得到

\[w_m=\frac3{(m+1)(m+3)}.\]

这个式子在\(m=0\)也等于1。利用

\[\frac1{(m+1)(m+3)}=\frac12\left(\frac1{m+1}-\frac1{m+3}\right),\]

可得

\[\sum_{m=0}^\infty w_m=\frac32\left(1+\frac12\right)=\frac94.\]

所以

\[\boxed{\pi_m=\frac4{3(m+1)(m+3)},\qquad m\ge0.}\]

所有项非负、总和1，相邻详细平衡已验证；非相邻双向转移均为0，故完整平稳方程成立。这才完成了``存在平稳概率''的证明，不能止于未归一化权重。

\subsection{变式1.72A：返回某一高度}\label{ux53d8ux5f0f1.72aux8fd4ux56deux67d0ux4e00ux9ad8ux5ea6}

\textbf{题目。}求从\(m\)出发第一次正返回\(m\)的平均时间。

\textbf{解答。}该不可约链已有平稳概率，故正常返。书内正返回公式给

\[\boxed{\mathbb E_mT_m=\frac1{\pi_m}=\frac{3(m+1)(m+3)}4.}\]

特别是从0返回的均值为\(9/4\)。初次原地停留也计一次正返回，所以答案不需要大于4。

\subsection{变式1.72B：正常返是否保证平稳状态的均值有限}\label{ux53d8ux5f0f1.72bux6b63ux5e38ux8fd4ux662fux5426ux4fddux8bc1ux5e73ux7a33ux72b6ux6001ux7684ux5747ux503cux6709ux9650}

\textbf{题目。}计算\(\mathbb E_\pi X\)是否有限。

\textbf{解答。}其级数为\(\sum_{m\ge1}4m/[3(m+1)(m+3)]\)。对\(m\ge3\)，有\(m+1\le2m,m+3\le2m\)，故

\[\frac{4m}{3(m+1)(m+3)}\ge\frac1{3m}.\]

调和级数发散，所以\(\boxed{\mathbb E_\pi X=\infty}\)。这不矛盾：正常返要求的是返回固定状态的平均时间有限，不要求任意状态函数在平衡下可积。

\subsection{变式1.72C：精确平稳尾概率}\label{ux53d8ux5f0f1.72cux7cbeux786eux5e73ux7a33ux5c3eux6982ux7387}

\textbf{题目。}求\(\mathbb P_\pi(X\ge M)\)，\(M\ge0\)。

\textbf{解答。}将原题的部分分式从\(M\)加到无穷，中间项抵消，得到

\[\boxed{\mathbb P_\pi(X\ge M)=\frac23\left(\frac1{M+1}+\frac1{M+2}\right).}\]

例如\(M=10\)时为\(23/198\)。\(M=0\)给1，\(M=1\)给\(5/9=1-\pi_0\)，同时检查了两端索引。

\section{习题1.73：为什么这一次不存在平稳概率}\label{ux4e60ux98981.73ux4e3aux4ec0ux4e48ux8fd9ux4e00ux6b21ux4e0dux5b58ux5728ux5e73ux7a33ux6982ux7387}

\textbf{书内依据：第1.4节平稳方程，第1.11节无穷状态空间。}状态从1开始，向上概率\(p_m=m/(2m+2)\)，\(m\ge1\)；向下概率\(q_m=1/2\)，\(m\ge2\)；内部停留概率\(1/(2m+2)\)，而\(P(1,1)=3/4\)。

\subsection{第一步：证明相邻平衡是必要的，不只是试猜}\label{ux7b2cux4e00ux6b65ux8bc1ux660eux76f8ux90bbux5e73ux8861ux662fux5fc5ux8981ux7684ux4e0dux53eaux662fux8bd5ux731c}

假设存在平稳概率\(\pi\)。取有限集合\(A_m=\{1,\ldots,m\}\)。平稳时一步之后在\(A_m\)的总概率不变，因此从集合流出的概率等于流入的概率。由于每步只在相邻状态之间移动，唯一向外的边是\(m\to m+1\)，唯一向内的边是\(m+1\to m\)。所以必有

\[\pi_m\frac{m}{2m+2}=\pi_{m+1}\frac12.\]

由此\(\pi_{m+1}=\pi_m m/(m+1)\)。连乘得到

\[\pi_m=\frac{\pi_1}{m}.\]

\subsection{第二步：归一化产生矛盾}\label{ux7b2cux4e8cux6b65ux5f52ux4e00ux5316ux4ea7ux751fux77dbux76fe}

若\(\pi_1>0\)，则\(\sum_m\pi_m=\pi_1\sum_m1/m=\infty\)，不是1。若\(\pi_1=0\)，递推又使所有\(\pi_m=0\)，总和为0。非负概率只能属于这两种情况，所以

\[\boxed{\text{本链不存在平稳概率分布。}}\]

不能只写``我的细致平衡尝试没成功，所以没有平稳分布''。上面的有限集合论证说明了任何平稳分布都被迫具有这个不可归一化的形式，才真正排除了所有可能。

\subsection{变式1.73A：没有平稳概率是否意味着暂留}\label{ux53d8ux5f0f1.73aux6ca1ux6709ux5e73ux7a33ux6982ux7387ux662fux5426ux610fux5473ux7740ux6682ux7559}

\textbf{题目。}进一步判断本链是否常返。

\textbf{解答。}将最低状态1当作1.70的边界。对\(m\ge2\)，有\(q_m/p_m=(m+1)/m\)，相邻首达差的乘积随高度线性增加，故常返判别中的正项级数发散。因此链常返。原题又证明不存在平稳概率，由不可约链判别，它是\(\boxed{\text{零常返}}\)，而不是暂留。

\subsection{变式1.73B：先截成有限链}\label{ux53d8ux5f0f1.73bux5148ux622aux6210ux6709ux9650ux94fe}

\textbf{题目。}截到\(1,\ldots,N\)，在\(N\)把向上概率改成原地概率，其他不变。求平稳分布并考察\(N\to\infty\)。

\textbf{解答。}边界改动不改变相邻边的双向概率，因此平稳权重仍为\(1/m\)。有限归一化后

\[\boxed{\pi_m^{(N)}=\frac1{mH_N},\quad1\le m\le N.}\]

返回1的平均时间为\(H_N\)，随\(N\)发散。对每个固定\(m\)，\(\pi_m^{(N)}\to0\)；这些点态极限之和为0，不是1。有限截断的平稳概率存在，并不意味着无穷链有平稳概率极限。

\subsection{变式1.73C：先到高端还是低端}\label{ux53d8ux5f0f1.73cux5148ux5230ux9ad8ux7aefux8fd8ux662fux4f4eux7aef}

\textbf{题目。}在首次到达1或\(N\)时停止，求从\(m\in\{1,\ldots,N\}\)先到\(N\)的概率。

\textbf{解答。}记\(h_1=0,h_N=1\)，令\(\Delta_m=h_m-h_{m-1}\)。内部差分为\(\Delta_{m+1}=((m+1)/m)\Delta_m\)，所以\(\Delta_j=(j/2)\Delta_2\)。相加得到\(h_m=\Delta_2\sum_{j=2}^m j/2=\Delta_2[m(m+1)-2]/4\)。以\(h_N=1\)归一化，得到

\[\boxed{h_m=\frac{m(m+1)-2}{N(N+1)-2},\qquad N\ge2.}\]

边界和每个内部方程都可代入检查。

\section{习题1.74：年龄增长、重置与返回}\label{ux4e60ux98981.74ux5e74ux9f84ux589eux957fux91cdux7f6eux4e0eux8fd4ux56de}

\textbf{书内依据：第1.7节返回时间，第1.11节正常返与平稳分布，非负整数变量的尾和公式。}在年龄\(n\)，以概率\(p_n\)增到\(n+1\)，以概率\(1-p_n\)重置为0，其中\(0\le p_n\le1\)。

定义

\[R_0=1,\qquad R_n=\prod_{k=0}^{n-1}p_k\quad(n\ge1).\]

从0出发，前\(n\)步尚未返回0，当且仅当连续\(n\)次都选择增长，因此

\[\mathbb P_0(T_0>n)=R_n.\]

\subsection{（a）常返的充要条件}\label{aux5e38ux8fd4ux7684ux5145ux8981ux6761ux4ef6}

事件\(\{T_0>n\}\)递减到\(\{T_0=\infty\}\)，所以

\[\boxed{0\text{常返}\iff\lim_{n\to\infty}R_n=0.}\]

每一步都有一定重置概率，不自动保证常返；这些概率若衰减得过快，仍可能永不重置。

\subsection{（b）正常返的充要条件}\label{bux6b63ux5e38ux8fd4ux7684ux5145ux8981ux6761ux4ef6}

非负整数值等待时间满足\(\mathbb ET=\sum_{n\ge0}\mathbb P(T>n)\)，即把一条长度\(T\)的路径数成\(T\)个单位。故

\[\mathbb E_0T_0=\sum_{n=0}^\infty R_n.\]

于是

\[\boxed{0\text{正常返}\iff\sum_{n=0}^\infty R_n<\infty.}\]

级数有限自动使其项趋于0，所以也自动包含了常返。

\subsection{（c）平稳分布}\label{cux5e73ux7a33ux5206ux5e03}

正状态\(n\)只有从\(n-1\)增长才能进入，因此任何平稳分布必满足\(\pi_n=\pi_{n-1}p_{n-1}\)，即\(\pi_n=\pi_0R_n\)。令\(S=\sum_{n\ge0}R_n\)。当\(S<\infty\)，归一化给

\[\boxed{\pi_n=R_n/S.}\]

还要核对0处的方程：\(R_n(1-p_n)=R_n-R_{n+1}\)，故

\[\sum_{n\ge0}\pi_n(1-p_n)=S^{-1}(R_0-\lim R_n)=1/S=\pi_0.\]

当\(S=\infty\)，\(\pi_n=\pi_0R_n\)无法归一化，所以没有平稳概率。这一结论也处理了某个\(p_n=0\)导致只能进入有限年龄类的情况；不应把不可达高年龄也一概说成常返状态。

\subsection{变式1.74A：每一步有重置机会但均值无穷}\label{ux53d8ux5f0f1.74aux6bcfux4e00ux6b65ux6709ux91cdux7f6eux673aux4f1aux4f46ux5747ux503cux65e0ux7a77}

\textbf{题目。}\(p_n=(n+1)/(n+2)\)，求分类。

\textbf{解答。}乘积抵消得\(R_n=1/(n+1)\)。它趋于0，所以0常返；但\(\sum R_n\)是调和级数，均值无穷。因此0零常返、无平稳概率。每次重置概率虽正，却不足以保证有限平均周期。

\subsection{变式1.74B：加大重置概率}\label{ux53d8ux5f0f1.74bux52a0ux5927ux91cdux7f6eux6982ux7387}

\textbf{题目。}\(p_n=(n+1)/(n+3)\)，求平稳分布。

\textbf{解答。}乘积为\(R_n=2/[(n+1)(n+2)]\)，总和为2。故

\[\boxed{\mathbb E_0T_0=2,\qquad\pi_n=\frac1{(n+1)(n+2)}.}\]

0正常返，且\(\pi_0=1/2\)与返回均值互为倒数。

\subsection{变式1.74C：永不重置的概率严格为正}\label{ux53d8ux5f0f1.74cux6c38ux4e0dux91cdux7f6eux7684ux6982ux7387ux4e25ux683cux4e3aux6b63}

\textbf{题目。}\(p_n=1-1/(n+2)^2\)，求从0永不正返回的概率。

\textbf{解答。}分解为\(p_n=(n+1)(n+3)/(n+2)^2\)，所以

\[R_n=\prod_{k=0}^{n-1}\frac{k+1}{k+2}\frac{k+3}{k+2}
=\frac{n+2}{2(n+1)}.\]

\(n=0\)给1，\(n=1\)给\(3/4\)，端点检查正确。令\(n\to\infty\)得到

\[\boxed{\mathbb P_0(T_0=\infty)=1/2.}\]

0暂留，无平稳概率。虽然每个有限年龄都能重置，仍有一半概率一直增长。

\section{习题1.75：先抽一个高度、再逐级下降}\label{ux4e60ux98981.75ux5148ux62bdux4e00ux4e2aux9ad8ux5ea6ux518dux9010ux7ea7ux4e0bux964d}

\textbf{书内依据：第1.7节一个返回周期，第1.11节正常返。}状态\(n>0\)下一步确定到\(n-1\)；在0处抽取一个有限非负整数\(J\)，分布\(\mathbb P(J=j)=p_j\)，其中\(\sum_{j\ge0}p_j=1\)，随后跳到\(J\)。

\subsection{原题：常返与正常返}\label{ux539fux9898ux5e38ux8fd4ux4e0eux6b63ux5e38ux8fd4}

从0跳到\(J\)用一步，之后确定下降\(J\)步回到0。因此逐条路径上都有

\[\boxed{T_0=J+1.}\]

由于\(J\)是正常概率分布下的有限整数，\(T_0<\infty\)的概率为1。因此0总是常返。进一步

\[\mathbb E_0T_0=1+\mathbb EJ=1+\sum_{j\ge0}jp_j,\]

所以

\[\boxed{0\text{正常返}\iff\sum_{j\ge0}jp_j<\infty.}\]

严格说，若\(J\)的支持有上界，则更高且从0不可达的状态可能是暂留的；``链常返''的说法应理解为0所在的不可约闭类常返。若支持无上界，所有有限状态都可从0经某个更高\(J\)到达，链不可约，才可以对全部状态统一这样说。

\subsection{补充：把一个周期转换成平稳分布}\label{ux8865ux5145ux628aux4e00ux4e2aux5468ux671fux8f6cux6362ux6210ux5e73ux7a33ux5206ux5e03}

若\(\mathbb EJ<\infty\)，从一次0到下一次0之间，在时刻\(0,\ldots,T_0-1\)内访问0恰一次，访问\(k\ge1\)恰在\(J\ge k\)时一次。故每周期在\(k\)的平均次数为\(\mathbb P(J\ge k)\)，平均周期为\(1+\mathbb EJ\)。于是

\[\boxed{\pi_k=\frac{\mathbb P(J\ge k)}{1+\mathbb EJ},\qquad k\ge0.}\]

尾和公式保证这些概率和为1；平稳方程还可直接核对为\(\pi_k=\pi_{k+1}+\pi_0p_k\)，因为\(\mathbb P(J\ge k)-\mathbb P(J\ge k+1)=p_k\)。

\subsection{变式1.75A：几何高度}\label{ux53d8ux5f0f1.75aux51e0ux4f55ux9ad8ux5ea6}

\textbf{题目。}\(p_j=(1-a)a^j\)，\(0\le a<1\)，求返回均值与平稳分布。

\textbf{解答。}\(\mathbb P(J\ge k)=a^k\)，\(\mathbb EJ=a/(1-a)\)。所以

\[\boxed{\mathbb E_0T_0=1/(1-a),\qquad\pi_k=(1-a)a^k.}\]

\(a=0\)时0吸收，公式变为\(\pi_0=1\)，其余为0。

\subsection{变式1.75B：下降总会结束，但平均时间无穷}\label{ux53d8ux5f0f1.75bux4e0bux964dux603bux4f1aux7ed3ux675fux4f46ux5e73ux5747ux65f6ux95f4ux65e0ux7a77}

\textbf{题目。}\(p_j=1/[(j+1)(j+2)]\)，求返回性质。

\textbf{解答。}\(p_j=1/(j+1)-1/(j+2)\)，所以总和1，是合法分布；尾概率\(\mathbb P(J\ge k)=1/(k+1)\)。从而\(\mathbb EJ=\sum_{k\ge1}1/(k+1)=\infty\)。每次抽出的\(J\)几乎必然有限，下降也必然结束，但平均周期无穷。因此链零常返，没有平稳概率。

\subsection{变式1.75C：正常返但有周期}\label{ux53d8ux5f0f1.75cux6b63ux5e38ux8fd4ux4f46ux6709ux5468ux671f}

\textbf{题目。}\(J=1,3\)各以概率\(1/2\)出现，求平稳分布、周期与逐步收敛情况。

\textbf{解答。}\(\mathbb EJ=2\)，平均周期3。尾概率依次为\(1,1,1/2,1/2\)，所以在\(0,1,2,3\)上

\[\boxed{\pi=(1/3,1/3,1/6,1/6).}\]

从0的第一次正返回时间只可能为2或4；任何多次返回长度都是偶数，又有长度2的返回，所以周期为2。从0出发，状态奇偶性每步互换，分布不逐步收敛到\(\pi\)。正常返、平稳分布存在和非周期性是三个不同条件。

\section{习题1.76：每家三个孩子的女性分枝过程}\label{ux4e60ux98981.76ux6bcfux5bb6ux4e09ux4e2aux5b69ux5b50ux7684ux5973ux6027ux5206ux679dux8fc7ux7a0b}

\textbf{书内依据：第1.11节分枝过程例1.55及最小不动点结论。}这里计算的是沿女性后代延续的家系。采用题设隐含的标准独立模型：每个孩子独立以概率\(1/2\)为女孩，各母亲独立生育；只说明平均有1.5个女儿，本身不足以唯一确定整个后代分布。

\subsection{第一步：女儿数分布与递推}\label{ux7b2cux4e00ux6b65ux5973ux513fux6570ux5206ux5e03ux4e0eux9012ux63a8}

一个母亲的女儿数\(K\)服从\(\operatorname{Bin}(3,1/2)\)，故概率生成函数为

\[f(s)=\mathbb E[s^K]=\frac18+\frac38s+\frac38s^2+\frac18s^3=\frac{(1+s)^3}{8}.\]

令\(q_n\)为从一位女性开始，到第\(n\)代已经灭绝的概率，\(q_0=0\)。第一代有\(K\)个女儿时，各分枝独立，下一阶段全部灭绝的概率为\(q_n^K\)。因此\(q_{n+1}=f(q_n)\)。

\(q_n\)递增且不超过1，故极限\(q\)存在，连续性给\(q=f(q)\)。若\(s\in[0,1]\)也是不动点，则从\(q_0=0\le s\)出发归纳得\(q_n\le s\)，所以\(q\le s\)。这证明必须选最小非负不动点，而不是见到根1就停止。

\subsection{第二步：因式分解并选对根}\label{ux7b2cux4e8cux6b65ux56e0ux5f0fux5206ux89e3ux5e76ux9009ux5bf9ux6839}

方程\(8q=(1+q)^3\)化为

\[q^3+3q^2-5q+1=(q-1)(q^2+4q-1)=0.\]

三个根为\(1,-2+\sqrt5,-2-\sqrt5\)。在\([0,1]\)中的最小根是

\[\boxed{q=\sqrt5-2\approx0.23606798.}\]

所以女性家系永续的概率为\(1-q=3-\sqrt5\)。这里不是说某一家``没有三个孩子''：总孩子数仍是三个，只是有可能没有女儿，或女儿的后续女性分枝都灭绝。

\subsection{变式1.76A：多位初始女性}\label{ux53d8ux5f0f1.76aux591aux4f4dux521dux59cbux5973ux6027}

\textbf{题目。}有\(k\)位相互独立的初始女性，求最终全部灭绝的概率。

\textbf{解答。}每条女性家系独立，且每条灭绝概率\(q=\sqrt5-2\)。全部灭绝是\(k\)个独立事件同时发生，所以

\[\boxed{q_k=(\sqrt5-2)^k.}\]

\(k=2\)时为\(9-4\sqrt5\)；至少一条延续的概率为\(1-q_k\)。这不是\(kq\)，因为这里要求的是``全部''，不是互斥事件的和。

\subsection{变式1.76B：第二代以内灭绝}\label{ux53d8ux5f0f1.76bux7b2cux4e8cux4ee3ux4ee5ux5185ux706dux7edd}

\textbf{题目。}求至第二代已无女性后代的概率。

\textbf{解答。}第一代灭绝概率\(q_1=f(0)=1/8\)。第二代只需每个女儿自己的下一代都没有女儿，所以

\[\boxed{q_2=f(q_1)=\frac{(1+1/8)^3}{8}=\frac{729}{4096}.}\]

它小于最终灭绝概率，因为还有第二代尚未灭绝、以后才灭绝的路径。不能把有限代灭绝概率当成最终概率。

\subsection{变式1.76C：女孩概率不再是二分之一}\label{ux53d8ux5f0f1.76cux5973ux5b69ux6982ux7387ux4e0dux518dux662fux4e8cux5206ux4e4bux4e00}

\textbf{题目。}每个孩子独立为女孩的概率是\(\theta\in[0,1]\)，每家仍三个孩子，求灭绝概率的临界条件及超临界公式。

\textbf{解答。}\(f(s)=(1-\theta+\theta s)^3\)。若\(3\theta\le1\)，则\(0\le f'(s)\le1\)，函数\(f(s)-s\)在\([0,1]\)不增、在1处为0；除退化端点的直接情形外，在\(s<1\)严格为正，所以唯一不动点为1，必然灭绝。

当\(\theta>1/3\)，\(f'(1)=3\theta>1\)，故\(s<1\)且充分接近1时\(f(s)-s<0\)；而\(f(0)\ge0\)。由连续性出现较小根。将\(f(s)-s\)除以\(s-1\)，所得二次因子为

\[\theta^3s^2+\theta^2(3-2\theta)s-(1-\theta)^3.\]

因此最小非负根为

\[\boxed{q=\frac{-\theta^2(3-2\theta)+\sqrt{\theta^4(3-2\theta)^2+4\theta^3(1-\theta)^3}}{2\theta^3}
\quad(\theta>1/3).}\]

\(\theta=1\)给\(q=0\)；\(\theta=1/2\)给原题的\(\sqrt5-2\)。根的选择依据仍是从0迭代得到的最小不动点。

\section{习题1.77：几何后代数的灭绝概率}\label{ux4e60ux98981.77ux51e0ux4f55ux540eux4ee3ux6570ux7684ux706dux7eddux6982ux7387}

\textbf{书内依据：第1.11节分枝过程、最小不动点；等比级数。}每个个体独立生育\(K\)个后代，\(\mathbb P(K=k)=p(1-p)^k\)，\(k\ge0\)，其中\(0<p\le1\)。这里的几何变量从0开始计，不是从1开始。

\subsection{第一步：先辨认参数约定}\label{ux7b2cux4e00ux6b65ux5148ux8fa8ux8ba4ux53c2ux6570ux7ea6ux5b9a}

生成函数为

\[f(s)=\sum_{k=0}^{\infty}p(1-p)^ks^k=\frac{p}{1-(1-p)s},\qquad0\le s\le1.\]

平均后代数是\(f'(1)=(1-p)/p\)。所以\(p\)越大，平均孩子数反而越少；不能把\(p\)当成生育成功概率来直觉判断增长。

\subsection{第二步：求不动点并分参数范围}\label{ux7b2cux4e8cux6b65ux6c42ux4e0dux52a8ux70b9ux5e76ux5206ux53c2ux6570ux8303ux56f4}

先令\(0<p<1\)。设灭绝概率为\(\rho\)，由\(\rho=f(\rho)\)得到

\[(1-p)\rho^2-\rho+p=(\rho-1)((1-p)\rho-p)=0.\]

两个候选根为\(1\)与\(p/(1-p)\)。取\([0,1]\)内最小根，得到

\[\boxed{\rho=
\begin{cases}
p/(1-p),&0<p<1/2,\\
1,&1/2\le p\le1.
\end{cases}}\]

\(p=1\)单独理解：所有个体都没有后代，一代后必灭绝。\(p=1/2\)为重根1，也必灭绝；不能因平均后代数恰为1而认为每代都保留一个个体。

\subsection{变式1.77A：有限代灭绝概率与临界等待时间}\label{ux53d8ux5f0f1.77aux6709ux9650ux4ee3ux706dux7eddux6982ux7387ux4e0eux4e34ux754cux7b49ux5f85ux65f6ux95f4}

\textbf{题目。}求\(q_n=\mathbb P(Z_n=0\mid Z_0=1)\)，并说明临界\(p=1/2\)时灭绝时间的均值。

\textbf{解答。}\(q_0=0,q_{n+1}=p/[1-(1-p)q_n]\)。令\(r=p/(1-p)\)。对\(0<p<1\)且\(p\ne1/2\)，直接代入递推并归纳可得

\[\boxed{q_n=\frac{r(1-r^n)}{1-r^{n+1}}.}\]

\(n=0\)为0；将第\(n\)项代入右边，分母通分后正好化成第\(n+1\)项，所以归纳闭合。即使\(r>1\)，分子分母同负，概率仍合法。

\(p=1/2\)时单独归纳得到\(q_n=n/(n+1)\)。若\(\tau=\min\{n\ge0:Z_n=0\}\)，则灭绝吸收性给\(\mathbb P(\tau>n)=1-q_n=1/(n+1)\)。故

\[\boxed{\mathbb P(\tau<\infty)=1,\qquad\mathbb E\tau=\infty\quad(p=1/2).}\]

又一次出现了``最终一定结束，但平均等待无穷''的区别。\(p=1\)时\(q_0=0\)、\(q_n=1\)对\(n\ge1\)成立，不套用\(r\)的表达式。

\subsection{变式1.77B：条件于最终灭绝，生育规律如何改变}\label{ux53d8ux5f0f1.77bux6761ux4ef6ux4e8eux6700ux7ec8ux706dux7eddux751fux80b2ux89c4ux5f8bux5982ux4f55ux6539ux53d8}

\textbf{题目。}在\(0<p<1/2\)时，已知最终家系灭绝，求第一代后代数的条件分布。

\textbf{解答。}原灭绝概率\(\rho=p/(1-p)\)。若第一代有\(k\)个后代，这\(k\)条独立家系都灭绝的概率为\(\rho^k\)。用条件概率定义，

\[\mathbb P(K=k\mid\text{最终灭绝})
=\frac{p(1-p)^k\rho^k}{\rho}
=\boxed{(1-p)p^k}.\]

条件化把几何参数由\(p\)换成\(1-p\)，均值由\((1-p)/p>1\)换成\(p/(1-p)<1\)。同样的分枝分解说明各子树在条件化后仍按此规律独立繁衍。包括祖先的条件总人口均值为等比和

\[\boxed{\mathbb E[\text{总人口}\mid\text{灭绝}]=\frac1{1-p/(1-p)}=\frac{1-p}{1-2p}.}\]

这不同于无条件总人口的期望，不能混合计算。

\subsection{变式1.77C：每代均值与总人口}\label{ux53d8ux5f0f1.77cux6bcfux4ee3ux5747ux503cux4e0eux603bux4ebaux53e3}

\textbf{题目。}从一人开始，求\(\mathbb EZ_n\)及包括祖先的总人口均值。

\textbf{解答。}给定\(Z_n=z\)，下一代是\(z\)个独立后代数之和，条件均值\(z\mu\)，其中\(\mu=(1-p)/p\)。因此\(\mathbb EZ_n=\mu^n\)。总人口\(S=\sum_{n\ge0}Z_n\)非负，先对有限代求和再令代数增长，期望单调上升，故

\[\boxed{\mathbb ES=\sum_{n\ge0}\mu^n=
\begin{cases}
p/(2p-1),&p>1/2,\\
\infty,&0<p\le1/2.
\end{cases}}\]

临界时虽然几乎必然灭绝、实际总人口几乎必然有限，期望仍可能无穷。\(p=1\)的总人口就是祖先一人，公式也给1。

\section{补齐第001批：习题1.1---1.7的21道变式}\label{ux8865ux9f50ux7b2c001ux6279ux4e60ux98981.11.7ux768421ux9053ux53d8ux5f0f}

这七道原题的解答已在第001批交付，本节不重复计算为新增原题，只补上每题三道不同训练。变式的内容和解答均为本项目自编，不冒充原书新题。对应原题分别是相邻掷币和、两盒交换球、累计骰子点数模六、住房状态、吸收随机游走、出租车和二日天气。

\subsection{变式1.1A：掷币有偏后，还会不会失去马尔可夫性}\label{ux53d8ux5f0f1.1aux63b7ux5e01ux6709ux504fux540eux8fd8ux4f1aux4e0dux4f1aux5931ux53bbux9a6cux5c14ux53efux592bux6027}

\textbf{题目。}\(Y_n\)独立且\(\mathbb P(Y_n=1)=p\)，令\(X_n=Y_{n-1}+Y_n\)。对\(0\le p\le1\)完整分类。

\textbf{解答。}若\(0<p<1\)，事件\(X_1=0,X_2=1\)迫使\((Y_0,Y_1,Y_2)=(0,0,1)\)；于是下一步\(X_3=2\)的概率为\(p\)。而\(X_1=2,X_2=1\)迫使前三项为\((1,1,0)\)，于是\(X_3\le1\)，到2概率为0。两个历史都有正概率、现在同为1、下一步分布却不同，故不是马尔可夫链。若\(p=0\)或1，整个过程分别恒为0或2，是退化马尔可夫链。因此只有这两个端点例外。

\subsection{变式1.1B：增大状态空间修复信息损失}\label{ux53d8ux5f0f1.1bux589eux5927ux72b6ux6001ux7a7aux95f4ux4feeux590dux4fe1ux606fux635fux5931}

\textbf{题目。}改记\(Z_n=(Y_{n-1},Y_n)\)，求转移矩阵和平稳分布。

\textbf{解答。}新的抛掷独立于过去，且\(Z_{n+1}\)的第一分量就是\(Z_n\)的第二分量，所以保留有序对已经足够。按\(00,01,10,11\)排列，

\[P=\begin{pmatrix}1-p&p&0&0\\0&0&1-p&p\\1-p&p&0&0\\0&0&1-p&p\end{pmatrix}.\]

相邻两个原始掷币独立，因此自然候选\(\pi=((1-p)^2,(1-p)p,p(1-p),p^2)\)。例如进入00的概率为\([(1-p)^2+p(1-p)](1-p)=(1-p)^2\)，其余三列同样检查，故\(\pi P=\pi\)。新链是马尔可夫链，但把00、11以外的两个状态合并成``和为1''就又丢掉了必要信息。

\subsection{变式1.1C：不重叠的两次之和}\label{ux53d8ux5f0f1.1cux4e0dux91cdux53e0ux7684ux4e24ux6b21ux4e4bux548c}

\textbf{题目。}\(W_n=Y_{2n}+Y_{2n+1}\)是否是马尔可夫链？

\textbf{解答。}每个\(W_n\)使用不同的两个独立掷币，故\(W_n\)本身相互独立，同分布为\(\operatorname{Bin}(2,p)\)。因此每行转移概率都相同：

\[\boxed{P(i,\cdot)=((1-p)^2,2p(1-p),p^2),\quad i=0,1,2.}\]

独立同分布序列当然满足马尔可夫性。原题失败的关键不在``取和''，而在相邻的和共享了一个未被完整保留的随机变量。

\subsection{变式1.2A：只求平均颜色数}\label{ux53d8ux5f0f1.2aux53eaux6c42ux5e73ux5747ux989cux8272ux6570}

\textbf{题目。}两盒各五球、总共五白五黑，当前左盒\(i\)个白球；求\(n\)次交换后的期望。

\textbf{解答。}条件增量的均值是

\[\frac{(5-i)^2-i^2}{25}=1-\frac{2i}{5}.\]

因此\(\mathbb E[X_{n+1}\mid X_n=i]=1+3i/5\)。令均值为\(\mu_n\)，其固定值为\(5/2\)，减去后得到公比\(3/5\)的等比递推。故

\[\boxed{\mathbb E_iX_n=\frac52+\left(\frac35\right)^n\left(i-\frac52\right).}\]

只求均值不必计算整个\(6\times6\)矩阵的\(n\)次幂。

\subsection{变式1.2B：三个球的完整小模型}\label{ux53d8ux5f0f1.2bux4e09ux4e2aux7403ux7684ux5b8cux6574ux5c0fux6a21ux578b}

\textbf{题目。}两盒各三球，总共三白三黑，求矩阵与平稳分布。

\textbf{解答。}交换后增、减1的概率分别为\((3-i)^2/9,i^2/9\)，不变概率\(2i(3-i)/9\)。按\(0,1,2,3\)排列，

\[P=\frac19\begin{pmatrix}0&9&0&0\\1&4&4&0\\0&4&4&1\\0&0&9&0\end{pmatrix}.\]

相邻平衡给比例\(\pi_1=9\pi_0\)、\(\pi_2=\pi_1\)、\(\pi_3=\pi_2/9\)，归一化得到

\[\boxed{\pi=(1,9,9,1)/20.}\]

边界0下一步必到1，边界3必到2；不能把边界误当成吸收状态。

\subsection{变式1.2C：短时返回概率与平均返回时间}\label{ux53d8ux5f0f1.2cux77edux65f6ux8fd4ux56deux6982ux7387ux4e0eux5e73ux5747ux8fd4ux56deux65f6ux95f4}

\textbf{题目。}原来的五球模型从0出发，求两步后回到0的概率和平均正返回时间。

\textbf{解答。}第一步必到1，第二步从1到0概率为\(1^2/25\)，所以\(P^2(0,0)=1/25\)。平稳分布由相邻平衡得到\(\pi_i=\binom5i^2/252\)，其中\(\sum_i\binom5i^2=252\)。故

\[\boxed{P^2(0,0)=1/25,\qquad\mathbb E_0T_0=252.}\]

\(25\)是某个两步概率的倒数，并非平均返回时间；倒数公式用的是平稳概率\(\pi_0\)。

\subsection{变式1.3A：模六累计和的两步概率}\label{ux53d8ux5f0f1.3aux6a21ux516dux7d2fux8ba1ux548cux7684ux4e24ux6b65ux6982ux7387}

\textbf{题目。}原题每轮两枚四面骰，增量余数概率\(q=(3,2,2,2,3,4)/16\)。从0出发，两轮后余数仍为0的概率是多少？

\textbf{解答。}第一轮余数为\(k\)后，第二轮必须为\(-k\bmod6\)。六种第一轮情形互斥，故

\[P^2(0,0)=\sum_{k=0}^5q_kq_{-k}
=\frac{9+8+6+4+6+8}{256}=\boxed{\frac{41}{256}}.\]

不是\(q_0^2=9/256\)，因为两轮的非零余数也可能相互抵消。

\subsection{变式1.3B：均匀平稳不代表不可约}\label{ux53d8ux5f0f1.3bux5747ux5300ux5e73ux7a33ux4e0dux4ee3ux8868ux4e0dux53efux7ea6}

\textbf{题目。}模六的增量改为0或2，各以\(1/2\)出现，求全部平稳分布。

\textbf{解答。}从\(i\)只能到\(i\)或\(i+2\pmod6\)，因此奇偶性保持。两个闭类为\(\{0,2,4\}\)和\(\{1,3,5\}\)。每个类内部形成带自环的三环，三点平稳概率相等。因此给偶类总质量\(\theta\)，给奇类\(1-\theta\)，得到全部平稳分布

\[\boxed{\pi=(\theta,1-\theta,\theta,1-\theta,\theta,1-\theta)/3,
\quad0\le\theta\le1.}\]

全六点均匀仅是\(\theta=1/2\)这一例；转移矩阵列和为1，并不能保证平稳分布唯一。

\subsection{变式1.3C：增量全是奇数}\label{ux53d8ux5f0f1.3cux589eux91cfux5168ux662fux5947ux6570}

\textbf{题目。}模六增量为1或3，各占一半。是否从每个状态收敛到均匀分布？

\textbf{解答。}每列之和仍为1，所以均匀分布平稳。反复取增量1可以走遍所有状态，链不可约。但每步都改变奇偶性，因此返回只在偶数步；同时两次取3就返回，总概率为\(1/4>0\)，所以周期恰为2。从0开始，奇数时刻永远不在0，而均匀平稳概率在0为\(1/6\)，不能逐步收敛。若再加入正概率增量0并保持增量1正概率，就有自环且仍不可约，才消除周期障碍。

\subsection{变式1.4A：任意起点的长期住房模型}\label{ux53d8ux5f0f1.4aux4efbux610fux8d77ux70b9ux7684ux957fux671fux4f4fux623fux6a21ux578b}

\textbf{题目。}原题十年转移规律不变，初始自有比例\(h_0\)任意，求第\(n\)期比例。

\textbf{解答。}\(h_{n+1}=0.12+0.82h_n\)，固定值满足\(h_*=0.12+0.82h_*\)，所以\(h_*=2/3\)。减去固定值并迭代，

\[\boxed{h_n=\frac23+\left(h_0-\frac23\right)\left(\frac{41}{50}\right)^n.}\]

这也是误差公式：若初始比例高于\(2/3\)便从上方下降，低于它便从下方上升。

\subsection{变式1.4B：长期比例不变是否意味着没有家庭迁移状态}\label{ux53d8ux5f0f1.4bux957fux671fux6bd4ux4f8bux4e0dux53d8ux662fux5426ux610fux5473ux7740ux6ca1ux6709ux5bb6ux5eadux8fc1ux79fbux72b6ux6001}

\textbf{题目。}在平衡下，每十年有多少比例的家庭改变拥有或租住状态？

\textbf{解答。}自有者占\(2/3\)，其中\(0.06\)转为租住，贡献\((2/3)0.06=0.04\)；租住者占\(1/3\)，其中\(0.12\)转为自有，贡献\((1/3)0.12=0.04\)。两方向互斥，相加得到\(\boxed{0.08=8\%}\)。净变化为0是两股流量相抵，不是流量本身都为0。

\subsection{变式1.4C：由长期比例与收敛速度反求概率}\label{ux53d8ux5f0f1.4cux7531ux957fux671fux6bd4ux4f8bux4e0eux6536ux655bux901fux5ea6ux53cdux6c42ux6982ux7387}

\textbf{题目。}另一二状态模型长期第一状态比例为\(0.6\)，偏离平衡的误差每期乘\(0.75\)。求转移矩阵。

\textbf{解答。}设第一状态转出概率为\(a\)、第二状态转入概率为\(b\)，则\(b/(a+b)=0.6\)，且\(1-a-b=0.75\)。所以\(a+b=0.25\)，\(b=0.15,a=0.10\)。矩阵为

\[\boxed{P=\begin{pmatrix}0.9&0.1\\0.15&0.85\end{pmatrix}.}\]

概率均在\([0,1]\)且每行和1，反向建模才算完成。

\subsection{变式1.5A：先赢到4还是先输到0}\label{ux53d8ux5f0f1.5aux5148ux8d62ux52304ux8fd8ux662fux5148ux8f93ux52300}

\textbf{题目。}原题向上概率\(2/5\)、向下\(3/5\)，0、4吸收，从1出发求先到4的概率。

\textbf{解答。}首达方程为\(2(h_{i+1}-h_i)=3(h_i-h_{i-1})\)，所以相邻差公比\(3/2\)。利用\(h_0=0,h_4=1\)，

\[\boxed{h_1=\frac1{1+3/2+(3/2)^2+(3/2)^3}=\frac8{65}.}\]

这个概率包括所有步数的获胜路径，不是前批计算的恰好三步后位于4的概率\(8/125\)。

\subsection{变式1.5B：平均多久退出}\label{ux53d8ux5f0f1.5bux5e73ux5747ux591aux4e45ux9000ux51fa}

\textbf{题目。}同一模型从1出发，求到0或4的平均时间。

\textbf{解答。}令\(g_0=g_4=0\)，则

\[g_1=1+\tfrac25g_2,\quad
g_2=1+\tfrac35g_1+\tfrac25g_3,\quad
g_3=1+\tfrac35g_2.\]

把第一、第三式代入第二式得\(g_2=2+(12/25)g_2\)，所以\(g_2=50/13\)。再代回，

\[\boxed{g_1=33/13,\qquad g_3=43/13.}\]

平均退出时间是\(33/13\)步。它不是成功所需的条件平均时间，两种退出结果都计入。

\subsection{变式1.5C：公平游戏的任意边界}\label{ux53d8ux5f0f1.5cux516cux5e73ux6e38ux620fux7684ux4efbux610fux8fb9ux754c}

\textbf{题目。}向上、向下各\(1/2\)，边界0与\(N\)吸收，初始为\(i\)，求获胜概率和平均时间。

\textbf{解答。}概率方程的相邻差恒定，边界给\(h_i=i/N\)。时间方程的二阶差为\(-2\)，而\(i(N-i)\)的二阶差正好是\(-2\)且两端为0。有限边界问题唯一，因此

\[\boxed{h_i=i/N,\qquad g_i=i(N-i).}\]

不要向含\(p-q\)分母的有偏公式直接代\(p=q\)，应像这里单独解退化后的差分方程。

\subsection{变式1.6A：出租车的长期位置}\label{ux53d8ux5f0f1.6aux51faux79dfux8f66ux7684ux957fux671fux4f4dux7f6e}

\textbf{题目。}机场\(A\)下一站等概率去酒店\(B,C\)；酒店以\(3/4\)去机场，以\(1/4\)去另一酒店。求平稳分布。

\textbf{解答。}对称性令酒店概率相等。设机场质量\(a\)，流入机场的总量为\((3/4)(1-a)\)，所以\(a=3(1-a)/4\)，得到\(a=3/7\)。剩余\(4/7\)由两酒店平分，

\[\boxed{\pi=(3/7,2/7,2/7).}\]

机场离开必去酒店，并不表示机场的长期质量为0。

\subsection{变式1.6B：返回机场的平均行程数}\label{ux53d8ux5f0f1.6bux8fd4ux56deux673aux573aux7684ux5e73ux5747ux884cux7a0bux6570}

\textbf{题目。}从机场出发，平均几次行程后正返回机场？

\textbf{解答。}无论当前在哪家酒店，下一次以\(3/4\)到机场，否则去另一酒店，所以从酒店到机场的剩余行程均值\(t\)满足\(t=1+t/4\)，得到\(t=4/3\)。从机场先用一次行程到酒店，因此

\[\boxed{\mathbb E_AT_A=1+4/3=7/3.}\]

也等于\(1/\pi_A\)，提供另一种检查。

\subsection{变式1.6C：第一次去某一家酒店}\label{ux53d8ux5f0f1.6cux7b2cux4e00ux6b21ux53bbux67d0ux4e00ux5bb6ux9152ux5e97}

\textbf{题目。}从机场出发，求第一次到\(B\)的平均行程数，允许途中回机场多次。

\textbf{解答。}设\(g_B=0\)，则\(g_A=1+g_C/2\)，\(g_C=1+3g_A/4\)。代入得\(g_A=3/2+3g_A/8\)，所以

\[\boxed{g_A=12/5,\qquad g_C=14/5.}\]

这与上一题返回机场的\(7/3\)不同，目标集合和起点不同，不能互相借用。

\subsection{变式1.7A：二日状态与单日下雨比例}\label{ux53d8ux5f0f1.7aux4e8cux65e5ux72b6ux6001ux4e0eux5355ux65e5ux4e0bux96e8ux6bd4ux4f8b}

\textbf{题目。}前两天都不雨时次日下雨概率\(0.3\)，否则为\(0.6\)。求按\(RR,RS,SR,SS\)排列的平稳分布，并求一天为雨天的长期比例。

\textbf{解答。}设四概率为\(a,b,c,d\)。转移规则给\(a=0.6a+0.6c\)、\(b=0.4a+0.4c\)、\(d=0.4b+0.7d\)，所以\(a=3c/2,b=c,d=4c/3\)。归一化后

\[\boxed{(a,b,c,d)=(9,6,6,8)/29.}\]

今天是否下雨看状态的第二字母，因此雨天比例是\(a+c=15/29\)，不是\(a+b\)的定义性解释；本例\(b=c\)让这两个数偶然相等，但应坚持正确的状态含义。

\subsection{变式1.7B：从连续两天不雨开始，多久第一次下雨}\label{ux53d8ux5f0f1.7bux4eceux8fdeux7eedux4e24ux5929ux4e0dux96e8ux5f00ux59cbux591aux4e45ux7b2cux4e00ux6b21ux4e0bux96e8}

\textbf{题目。}当前状态\(SS\)，求未来首次雨天的等待分布和均值。

\textbf{解答。}只要还不下雨，二日状态始终是\(SS\)，每个下一天的条件雨概率始终为\(0.3\)。因此等待\(W\)满足

\[\boxed{\mathbb P(W=k)=0.7^{k-1}0.3\ (k\ge1),\qquad\mathbb EW=10/3.}\]

即使整个天气序列不独立，这一特定等待阶段仍有几何分布，因为尚未成功时状态不变。

\subsection{变式1.7C：等到连续两天雨}\label{ux53d8ux5f0f1.7cux7b49ux5230ux8fdeux7eedux4e24ux5929ux96e8}

\textbf{题目。}从\(SS\)开始，求首次到\(RR\)的平均天数。

\textbf{解答。}记三个非目标状态的均值为\(g_{SS},g_{SR},g_{RS}\)，目标\(g_{RR}=0\)。第一步方程是

\[
\begin{aligned}
g_{SS}&=1+0.7g_{SS}+0.3g_{SR},\\
g_{SR}&=1+0.4g_{RS},\\
g_{RS}&=1+0.6g_{SR}+0.4g_{SS}.
\end{aligned}
\]

首式给\(g_{SS}=10/3+g_{SR}\)。代入末式得\(g_{RS}=7/3+g_{SR}\)；再代入中式，\(g_{SR}=29/15+0.4g_{SR}\)，所以\(g_{SR}=29/9\)。最终

\[\boxed{g_{SS}=59/9,\qquad g_{RS}=50/9.}\]

第一次下雨并不等于已经出现连续两天雨，中途失败后也不总是回到同一个状态，所以不能把本题当作成功概率\(0.3^2\)的几何等待。

\section{完成清单、范围与下一接续点}\label{ux5b8cux6210ux6e05ux5355ux8303ux56f4ux4e0eux4e0bux4e00ux63a5ux7eedux70b9}

\subsection{本批原题覆盖}\label{ux672cux6279ux539fux9898ux8986ux76d6}

本批新增原题题号连续为1.56---1.77，共22个，每个题目的全部小问均单列或在完整证明中处理。原题1.56---1.60分别涉及贷款吸收、五状态链、两状态递推、一般交换球和移至最前书架；1.61---1.69涉及圆环、棋盘、遗传、收集与随机下降；1.70---1.77完成无穷状态链和分枝过程。

每个新增原题有A、B、C三道变式，共66道；另补1.1---1.7每题A、B、C，共21道。因此本批新增自编训练87道，训练不计入原书题数。

\subsection{第1章累计账目}\label{ux7b2c1ux7ae0ux7d2fux8ba1ux8d26ux76ee}

\begin{longtable}[]{@{}llrr@{}}
\toprule\noalign{}
批次 & 原题范围 & 原题题号数 & 该批新增变式 \\
\midrule\noalign{}
\endhead
\bottomrule\noalign{}
\endlastfoot
001 & 1.1---1.7 & 7 & 0（其21道变式在004批补齐） \\
002 & 1.8---1.25 & 18 & 54 \\
003 & 1.26---1.55 & 30 & 90 \\
004 & 1.56---1.77 & 22 & 87（其中21道补前批） \\
合计 & 1.1---1.77 & 77 & 231 \\
\end{longtable}

至本批结束，Durrett第三版第1章的77个原题题号已连续处理，每个题号现均配齐三道自编详解变式，共231道。题面异常单独保留：例如1.9(c)不合法矩阵的原样无解证明已在第002批给出，不把自选修正版冒充原题；本批1.64、1.70、1.71等的适用条件也明确说明。

原题存在排印或条件边界问题时，``处理完成''表示已给出原样的正确判断及必要的条件说明，不表示为一个错误命题伪造肯定证明。

\subsection{五书状态}\label{ux4e94ux4e66ux72b6ux6001}

Durrett：第1章习题部分已经完成；下一道是\textbf{2016年第三版第2章习题2.1}。2021公开稿第2章已插入或改动部分习题，后续继续前必须核对版本，不能把同题号自动视为同题。

Ross与三本中文教材：本批没有新增原题解答；其余书的全量习题原文与题号清单仍需分别取得或完整核对。不能把完成Durrett一章写成五本书全部完成。

\subsection{计算检查与证明的区别}\label{ux8ba1ux7b97ux68c0ux67e5ux4e0eux8bc1ux660eux7684ux533aux522b}

资料包中的Python脚本使用有理数、符号运算或精确整数计数，检查有限矩阵行和、平稳关系、线性方程残差、棋盘度数、若干通式及变式数值。执行记录列出实际检查数量。无穷级数判别、条件独立、边界例外与一般参数结论依靠本讲义的文字证明；程序检查不是对全部数学证明的形式化认证。

\subsection{参考定位}\label{ux53c2ux8003ux5b9aux4f4d}

主要教材：Richard Durrett, \emph{Essentials of Stochastic Processes},
third edition, Springer, 2016, ISBN 9783319456133, DOI
10.1007/978-3-319-45614-0。第1章练习1.56---1.77位于2016版印刷页90---94。

交叉核对：作者公开稿 \emph{Essentials of Stochastic Processes}, Version
3.9, May
2021，Duke作者页面提供的PDF，对应印刷页73---77。本讲义的解答与变式为独立推导，只保留求解所需条件与必要矩阵，不复制全章英文题干。

作者公开稿：\url{https://sites.math.duke.edu/~rtd/EOSP/EOSP2021.pdf}。

\end{document}